% Lesson 4 — Speaking: Exchanges information about applying for a passport — Anglais Tle A — Cours signé OnBuch+
% Programme officiel MINESEC. Compiler avec : tectonic EN04-speaking-exchanges-information-about-applying-for-a-pas.tex
\newcommand{\DOCMATIERE}{Anglais}
\newcommand{\DOCNIVEAU}{Tle A}
\newcommand{\DOCMODULE}{Chapter 1 : Family and Social Life}
\newcommand{\DOCLECON}{EN04}
\newcommand{\DOCTITRE}{Lesson 4 — Speaking: Exchanges information about applying for a passport}
% OnBuch+ — préambule « cours signés » ENRICHI (compilé avec tectonic / XeLaTeX)
% Système visuel « Pop » d'OnBuch+ : fond crème (jamais blanc), bordures encre
% épaisses, ombres « dures » décalées, titres Archivo Black en capitales.
% Version MATHÉMATIQUES : graphes (pgfplots/tikz), boîtes pédagogiques
% (définition, propriété, méthode, attention, le savais-tu, exercice résolu,
% exercice, TP, bilan), page de garde et signature OnBuch+ ancrée en bas.
%
% Macros attendues dans main.tex AVANT \input{preamble} :
%   \DOCMATIERE  \DOCNIVEAU  \DOCMODULE  \DOCLECON (numéro)  \DOCTITRE
\documentclass[11pt]{article}

\usepackage{fontspec}
\usepackage[english,french]{babel} % français = langue principale ; l'anglais via \eng{..} / textebox
\usepackage[a4paper,margin=2cm,top=3.4cm,bottom=2.8cm,headheight=1.4cm,footskip=1.3cm]{geometry}
\usepackage{amsmath,amssymb,amsfonts}
\usepackage{mathtools} % \xmapsto, \coloneqq... souvent utilisés par les modèles
\usepackage{cancel} % simplifications barrées (bilans de forces)
% \abs{z} (module, valeur absolue) et \norm{u} (norme), utilisés par les modèles
\providecommand{\abs}{}\let\abs\relax\DeclarePairedDelimiter{\abs}{\lvert}{\rvert}
\providecommand{\norm}{}\let\norm\relax\DeclarePairedDelimiter{\norm}{\lVert}{\rVert}
\usepackage[table]{xcolor} % [table] : fournit aussi \rowcolors (alternance de lignes)
\usepackage{pagecolor}
\usepackage{tikz}
\usetikzlibrary{arrows.meta,positioning,calc,shapes.geometric,shapes.misc,shapes.arrows,decorations.pathmorphing,decorations.pathreplacing,decorations.markings,patterns,shadows,mindmap,trees,fit,backgrounds,matrix,intersections,angles,quotes}
\usepackage{pgfplots}
\usepackage[version=4]{mhchem} % équations nucléaires \ce{^{235}_{92}U}
\usepackage[european,siunitx]{circuitikz} % schémas électriques
\pgfplotsset{compat=1.18}
\usepgfplotslibrary{fillbetween,groupplots} % aires entre courbes (calcul intégral)
\usepackage{enumitem}
\usepackage{fancyhdr}
% [most] charge toutes les bibliothèques tcolorbox (listings, minted, raster,
% documentation...) et gonfle inutilement la mémoire TeX sur un document long
% (~300 boîtes) : on ne charge que ce qui est réellement utilisé.
\usepackage[skins,breakable]{tcolorbox}
\usepackage{graphicx}
\usepackage{colortbl}
\usepackage{booktabs}
\usepackage{tabularx}
\usepackage{diagbox} % case d'en-tête diagonale des tableaux à double entrée
% Colonnes extensibles centrée / gauche / droite (C, L, R), souvent utilisées par les modèles
\newcolumntype{C}{>{\centering\arraybackslash}X}
\newcolumntype{L}{>{\raggedright\arraybackslash}X}
\newcolumntype{R}{>{\raggedleft\arraybackslash}X}
\usepackage{array}
\usepackage{multicol}
\usepackage{multirow}
\usepackage{siunitx}
% parse-numbers=false : accepte aussi \SI{2,5 \times 10^{-3}}{...}
\sisetup{locale=FR,output-decimal-marker={,},group-minimum-digits=4,parse-numbers=false}
\DeclareSIUnit\ohm{\ensuremath{\symup{\Omega}}} % U+2126 absent de la police
\DeclareSIUnit\division{div} % réglages d'oscilloscope (ms/div, V/div)
\DeclareSIUnit\tr{tr} % tours (vitesse de rotation, ex. \SI{1500}{\tr\per\minute})
\DeclareSIUnit\molar{M} % concentration molaire (ex. \SI{3}{\molar}, \SI{1}{\micro\molar})
\DeclareSIUnit\an{an} % année (le modèle confond parfois avec \year, un compteur TeX) — ex. \SI{5}{\kilo\meter\per\an}
\DeclareSIUnit\FCFA{FCFA} % franc CFA (ex. \SI{500}{\FCFA\per\kilo\gram})
\DeclareSIUnit\degreeBrix{\textdegree Brix} % teneur en sucre d'un sirop/jus (réfractométrie)
\DeclareSIUnit\month{mois} % le modèle utilise parfois \month, un compteur TeX, comme unité de durée
\usepackage{qrcode}
% \degree : commande fréquemment utilisée par les modèles pour un angle ou une
% température en dehors de \SI{}{} (non fournie par les paquets chargés).
\providecommand{\degree}{^{\circ}}
% \qed (fin de démonstration), utilisé par les modèles sans amsthm
\providecommand{\qed}{\hfill$\square$}
% \barre : double barre des valeurs interdites dans un tableau de variations (tkz-tab)
\providecommand{\barre}{\ensuremath{\|}}
% environnement proof (amsthm non chargé) utilisé par les modèles
\ifdefined\proof\else\newenvironment{proof}[1][Démonstration]{\par\noindent\textit{#1.}\ }{\qed\par}\fi
\usepackage{lastpage}
\usepackage{hyperref}
\hypersetup{hidelinks}

% ============================================================
% Palette « Pop » OnBuch+ — mêmes hex que la classe P (plus_ui.dart)
% ============================================================
\definecolor{poporange}{HTML}{F59321}
\definecolor{popdark}{HTML}{E07A0C}
\definecolor{popcream}{HTML}{FFF6E8}
\definecolor{popink}{HTML}{1C1714}
\definecolor{poppurple}{HTML}{7A5AE0}
\definecolor{popgreen}{HTML}{2E9E6A}
\definecolor{poppink}{HTML}{E0457B}
\definecolor{popblue}{HTML}{2F7DE1}
\definecolor{popgold}{HTML}{C98A12}
\definecolor{popmuted}{HTML}{8A7F76}
\definecolor{popline}{HTML}{EADFD2}
\definecolor{popgray}{HTML}{8A7F76}
\colorlet{popgrey}{popgray}
% Alias tolérants (le modèle nomme parfois les couleurs différemment)
\colorlet{popwhite}{white}
\colorlet{popblack}{popink}
\colorlet{popred}{poppink}
\colorlet{popyellow}{popgold}
% Teintes claires (fonds de tableaux/encadrés) — le modèle nomme parfois
% les couleurs avec un suffixe L (ex. popblueL) sans les avoir définies.
\colorlet{poporangeL}{poporange!20}
\colorlet{popdarkL}{popdark!20}
\colorlet{poppurpleL}{poppurple!20}
\colorlet{popgreenL}{popgreen!20}
\colorlet{poppinkL}{poppink!20}
\colorlet{popblueL}{popblue!20}
\colorlet{popgoldL}{popgold!20}
\colorlet{popredL}{popred!20}
\colorlet{popgrayL}{popgray!20}
% Teintes claires pour fonds de tableaux / zones
\colorlet{popblueL}{popblue!10!white}
\colorlet{poporangeL}{poporange!14!white}
\colorlet{popgreenL}{popgreen!12!white}
\colorlet{poppurpleL}{poppurple!12!white}
\colorlet{poppinkL}{poppink!12!white}
\colorlet{popgoldL}{popgold!14!white}

\pagecolor{popcream}

% ============================================================
% Typographie
% ============================================================
\defaultfontfeatures{Ligatures=TeX}
\setmainfont{PlusJakartaSans-Medium.ttf}[Path=../../../fonts/,BoldFont=PlusJakartaSans-Bold.ttf]
% Symboles, API et emojis absents de la police principale : repli sur DejaVu Sans (caractères actifs)
\newfontface\symfont{DejaVuSans.ttf}[Path=../../../fonts/]
\makeatletter
\newcommand\symactive[1]{\catcode#1=13 \begingroup\lccode`\~=#1 \lowercase{\endgroup\def~}{{\symfont\char#1}}}
\newcommand\symblank[1]{\catcode#1=13 \begingroup\lccode`\~=#1 \lowercase{\endgroup\def~}{}}
\newcommand\symhyph[1]{\catcode#1=13 \begingroup\lccode`\~=#1 \lowercase{\endgroup\def~}{\mbox{-}}}
\newcount\symc
\newcommand\symrange[2]{\symc=#1 \@whilenum\symc<#2 \do{\expandafter\symactive\expandafter{\the\symc}\advance\symc by 1 }}
\newcommand\symblankrange[2]{\symc=#1 \@whilenum\symc<#2 \do{\expandafter\symblank\expandafter{\the\symc}\advance\symc by 1 }}
\makeatother
\symrange{"0250}{"02FF}\symactive{"2713}\symactive{"2714}\symactive{"2717}\symactive{"2718}\symactive{"2610}\symactive{"2611}\symactive{"2612}
\symhyph{"2011}\symhyph{"2E1A}\symblankrange{"1F300}{"1FAFF}\symblank{"FE0F}
% Maths en sans-serif (Fira Math) avec chiffres et lettres droites de Plus
% Jakarta, pour que les formules se fondent dans le texte.
\usepackage[math-style=ISO,bold-style=ISO]{unicode-math}
\setmathfont{FiraMath-Regular.otf}[Path=../../../fonts/]
\setmathfont{PlusJakartaSans-Medium.ttf}[Path=../../../fonts/,range={up/num,up/{Latin,latin},\mathup}]
\usepackage{icomma} % virgule décimale sans espace en mode math
% Caractères mathématiques Unicode absents des polices de texte (Plus Jakarta,
% Archivo Black) : rendus par leur équivalent mathématique, en texte comme en
% formule — sinon ils s'affichent en carré vide (ex. « Arithmétique dans ℤ »).
\usepackage{newunicodechar}
\newunicodechar{ℝ}{\ensuremath{\mathbb{R}}}
\newunicodechar{ℤ}{\ensuremath{\mathbb{Z}}}
\newunicodechar{ℂ}{\ensuremath{\mathbb{C}}}
\newunicodechar{ℕ}{\ensuremath{\mathbb{N}}}
\newunicodechar{ℚ}{\ensuremath{\mathbb{Q}}}
\newunicodechar{𝔻}{\ensuremath{\mathbb{D}}}
\newunicodechar{α}{\ensuremath{\alpha}}
\newunicodechar{β}{\ensuremath{\beta}}
\newunicodechar{γ}{\ensuremath{\gamma}}
\newunicodechar{δ}{\ensuremath{\delta}}
\newunicodechar{ε}{\ensuremath{\varepsilon}}
\newunicodechar{θ}{\ensuremath{\theta}}
\newunicodechar{λ}{\ensuremath{\lambda}}
\newunicodechar{μ}{\ensuremath{\mu}}
\newunicodechar{σ}{\ensuremath{\sigma}}
\newunicodechar{τ}{\ensuremath{\tau}}
\newunicodechar{φ}{\ensuremath{\varphi}}
\newunicodechar{ψ}{\ensuremath{\psi}}
\newunicodechar{ω}{\ensuremath{\omega}}
\newunicodechar{Σ}{\ensuremath{\Sigma}}
% majuscules produites par \MakeUppercase dans les titres (α-aminés -> Α)
\newunicodechar{Α}{\ensuremath{\alpha}}
\newunicodechar{Β}{\ensuremath{\beta}}
\newunicodechar{Γ}{\ensuremath{\Gamma}}
\newunicodechar{Δ}{\ensuremath{\Delta}}
\newunicodechar{Ε}{\ensuremath{\varepsilon}}
\newunicodechar{Θ}{\ensuremath{\Theta}}
\newunicodechar{Λ}{\ensuremath{\Lambda}}
\newunicodechar{Μ}{\ensuremath{\mu}}
\newunicodechar{Τ}{\ensuremath{\tau}}
\newunicodechar{Φ}{\ensuremath{\Phi}}
\newunicodechar{Ψ}{\ensuremath{\Psi}}
\newunicodechar{Ω}{\ensuremath{\Omega}}
\newunicodechar{↦}{\ensuremath{\mapsto}}
\newunicodechar{∈}{\ensuremath{\in}}
\newunicodechar{∉}{\ensuremath{\notin}}
\newunicodechar{∧}{\ensuremath{\wedge}}
\newunicodechar{⇔}{\ensuremath{\Leftrightarrow}}
\newunicodechar{⟺}{\ensuremath{\Longleftrightarrow}}
\newunicodechar{⇒}{\ensuremath{\Rightarrow}}
\newunicodechar{⟹}{\ensuremath{\Longrightarrow}}
\newunicodechar{∘}{\ensuremath{\circ}}
\newunicodechar{∪}{\ensuremath{\cup}}
\newunicodechar{∩}{\ensuremath{\cap}}
\newunicodechar{≡}{\ensuremath{\equiv}}
\newunicodechar{⊂}{\ensuremath{\subset}}
\newunicodechar{⊥}{\ensuremath{\perp}}
\newunicodechar{∃}{\ensuremath{\exists}}
\newunicodechar{∀}{\ensuremath{\forall}}
\newunicodechar{∛}{\ensuremath{\sqrt[3]{\,}}}
\newunicodechar{∜}{\ensuremath{\sqrt[4]{\,}}}
\newunicodechar{⃗}{\ensuremath{{}^{\to}}}
\newunicodechar{ⁿ}{\ifmmode^{n}\else\textsuperscript{\ensuremath{n}}\fi}
\newunicodechar{ˣ}{\ifmmode^{x}\else\textsuperscript{\ensuremath{x}}\fi}
\newunicodechar{ᵇ}{\ifmmode^{b}\else\textsuperscript{\ensuremath{b}}\fi}
\newunicodechar{ʳ}{\ifmmode^{r}\else\textsuperscript{\ensuremath{r}}\fi}
\newunicodechar{ᵘ}{\ifmmode^{u}\else\textsuperscript{\ensuremath{u}}\fi}
\newunicodechar{ʸ}{\ifmmode^{y}\else\textsuperscript{\ensuremath{y}}\fi}
\newunicodechar{ᵃ}{\ifmmode^{a}\else\textsuperscript{\ensuremath{a}}\fi}
\newunicodechar{ᵉ}{\ifmmode^{e}\else\textsuperscript{\ensuremath{e}}\fi}
\newunicodechar{ᵏ}{\ifmmode^{k}\else\textsuperscript{\ensuremath{k}}\fi}
\newunicodechar{ᵖ}{\ifmmode^{p}\else\textsuperscript{\ensuremath{p}}\fi}
\newunicodechar{ᴺ}{\ifmmode^{N}\else\textsuperscript{\ensuremath{N}}\fi}
\newunicodechar{ᶜ}{\ifmmode^{c}\else\textsuperscript{\ensuremath{c}}\fi}
\newunicodechar{ʰ}{\ifmmode^{h}\else\textsuperscript{\ensuremath{h}}\fi}
\newunicodechar{ᵠ}{\ifmmode^{q}\else\textsuperscript{\ensuremath{q}}\fi}
\newunicodechar{ₙ}{\ifmmode_{n}\else\textsubscript{\ensuremath{n}}\fi}
\newunicodechar{ₐ}{\ifmmode_{a}\else\textsubscript{\ensuremath{a}}\fi}
\newunicodechar{ᵢ}{\ifmmode_{i}\else\textsubscript{\ensuremath{i}}\fi}
\newunicodechar{ᵣ}{\ifmmode_{r}\else\textsubscript{\ensuremath{r}}\fi}
\newunicodechar{ₖ}{\ifmmode_{k}\else\textsubscript{\ensuremath{k}}\fi}
\newunicodechar{ᵦ}{\ifmmode_{\beta}\else\textsubscript{\ensuremath{\beta}}\fi}
\newunicodechar{ₓ}{\ifmmode_{x}\else\textsubscript{\ensuremath{x}}\fi}
\newfontfamily\popdisplay{ArchivoBlack-Regular.ttf}[Path=../../../fonts/]
\newfontfamily\poplogo{SpaceGrotesk-Bold.ttf}[Path=../../../fonts/]
\newfontfamily\popsemi{PlusJakartaSans-SemiBold.ttf}[Path=../../../fonts/]
\newfontfamily\popxbold{PlusJakartaSans-ExtraBold.ttf}[Path=../../../fonts/]
\newfontfamily\popxboldls{PlusJakartaSans-ExtraBold.ttf}[Path=../../../fonts/,LetterSpace=3.0]

\setlist[enumerate]{leftmargin=1.7em,itemsep=5pt,topsep=4pt}
\setlist[itemize]{leftmargin=1.7em,itemsep=4pt,topsep=4pt,label={\color{poporange}\textbullet}}
\setlength{\parindent}{0pt}
\setlength{\parskip}{5pt}
\renewcommand{\arraystretch}{1.25}

% pgfplots : style Pop par défaut
\pgfplotsset{
  every axis/.append style={axis line style={popink,line width=1pt},tick style={popink},
    grid style={popline,line width=0.5pt},label style={font=\small},tick label style={font=\footnotesize}},
  popcurve/.style={poporange,line width=1.8pt,no markers,smooth},
}
% Même style disponible dans un \draw TikZ simple (hors pgfplots)
\tikzset{popcurve/.style={poporange,line width=1.8pt}}
\tikzset{popgrid/.style={popline,very thin}}
\colorlet{popcurve}{poporange} % le nom du style est parfois utilisé comme couleur

% ============================================================
% Pill / squiggle
% ============================================================
\newcommand{\poppill}[3]{%
  \tikz[baseline=-0.55ex]\node[draw=popink,line width=1.2pt,rounded corners=2.6pt,%
    fill=#2,inner xsep=6.5pt,inner ysep=3pt]{%
    \popxboldls\fontsize{7}{7}\selectfont\color{#3}\MakeUppercase{#1}};%
}
\newcommand{\popsquiggle}[1]{%
  \tikz[baseline=-0.4ex]{
    \draw[#1,line width=2.6pt,line cap=round]
      plot [smooth] coordinates {(0,0.09) (0.34,0.01) (0.68,0.21) (1.02,0.01) (1.36,0.10)};
  }%
}
% Mot-clé surligné (marqueur orange sous le texte)
\newcommand{\cle}[1]{{\popxbold\color{popdark}#1}}

% ============================================================
% En-tête / pied
% ============================================================
\pagestyle{fancy}
\fancyhf{}
\renewcommand{\headrulewidth}{0pt}
\renewcommand{\footrulewidth}{0pt}
\lhead{\raisebox{-0.15ex}{\poplogo\fontsize{13}{13}\selectfont\color{popink}OnBuch\color{poporange}+}}
\rhead{\poppill{\DOCMATIERE\ \textbullet\ \DOCNIVEAU\ \textbullet\ Leçon \DOCLECON}{white}{popink}}
\lfoot{\popxbold\fontsize{7.6}{7.6}\selectfont\color{popmuted}\MakeUppercase{Cours signé OnBuch+}\ \textbullet\ {\popsemi\DOCTITRE}}
\rfoot{\tikz[baseline=-0.6ex]\node[rounded corners=8.5pt,fill=popink,minimum height=17pt,minimum width=17pt,inner xsep=5pt,inner sep=0]{\popxbold\fontsize{8}{8}\selectfont\color{popcream}\thepage\,/\,\pageref*{LastPage}};}

% ============================================================
% Titres
% ============================================================
\newcommand{\chaptitre}[2]{%
  \begin{tcolorbox}[enhanced,colback=poporange,colframe=popink,boxrule=2.6pt,arc=11pt,
      drop shadow=popink,left=15pt,right=15pt,top=12pt,bottom=11pt,before skip=3pt,after skip=18pt]
    \poppill{#2}{popink}{popcream}\par\vspace{9pt}
    {\raggedright\hyphenpenalty=10000\exhyphenpenalty=10000\popdisplay\fontsize{22}{24}\selectfont\color{popcream}\MakeUppercase{#1}\par}
  \end{tcolorbox}
}
% Section numérotée : pastille orange avec le numéro + titre Archivo
\newcounter{popsec}
\newcommand{\coursec}[1]{%
  \stepcounter{popsec}\setcounter{popsub}{0}%
  \par\vspace{16pt}\noindent
  \begin{minipage}{\linewidth}
  \tikz[baseline=-0.7ex]\node[draw=popink,line width=1.6pt,fill=poporange,rounded corners=4pt,minimum size=22pt,inner sep=0]{\popdisplay\fontsize{12}{12}\selectfont\color{popcream}\thepopsec};%
  \hspace{8pt}{\popdisplay\fontsize{15}{17}\selectfont\color{popink}\MakeUppercase{#1}}\par
  \vspace{4pt}\noindent{\color{poporange}\rule{36pt}{3.6pt}}
  \end{minipage}\par\nopagebreak\vspace{8pt}
}
\newcounter{popsub}
\newcommand{\courssub}[1]{%
  \stepcounter{popsub}%
  \par\vspace{9pt}\noindent{\popxbold\fontsize{11.5}{13}\selectfont\color{popink}{\color{poporange}\thepopsec.\thepopsub}\hspace{6pt}#1}\par\nopagebreak\vspace{4pt}
}

% ============================================================
% Boîtes pédagogiques — même recette : fond blanc, bordure épaisse colorée,
% ombre dure encre, badge plein. Titre optionnel [#1].
% ============================================================
\tcbset{popbase/.style={breakable,enhanced,colback=white,boxrule=2.3pt,arc=9pt,
  shadow={3pt}{-3pt}{0pt}{popink},left=14pt,right=14pt,top=11pt,bottom=11pt,before skip=11pt,after skip=15pt,
  fontupper=\popsemi\fontsize{10.6}{14.5}\selectfont\color{popink},
  fontlower=\popsemi\fontsize{10.6}{14.5}\selectfont\color{popink}}}
% Coche dessinée (le glyphe ✓ n'existe pas dans les polices du document)
\newcommand{\popcheck}[1]{\tikz[baseline=-0.5ex]\draw[#1,line width=1.6pt,line cap=round,line join=round](-0.13,0.0)--(-0.04,-0.09)--(0.13,0.1);}
\providecommand{\coche}{\popcheck{popgreen}}
\providecommand{\resultat}[1]{\par\smallskip\noindent\fcolorbox{popgreen}{popgreenL}{\parbox{\dimexpr\linewidth-2\fboxsep-2\fboxrule\relax}{\textbf{Résultat.} #1}}\par\smallskip}
\newcommand{\popboxhead}[3]{\poppill{#1}{#2}{white}\ifx\relax#3\relax\else\hspace{6pt}{\popxbold\fontsize{10.6}{12}\selectfont\color{popink}#3}\fi\par\vspace{7pt}}

\newtcolorbox{aretenir}[1][]{popbase,colframe=popgreen,before upper={\popboxhead{À retenir}{popgreen}{#1}}}
% Texte en anglais (document de lecture, dialogue, extrait) : cadre
% d'étude. Un titre optionnel (source, auteur) s'affiche dans l'en-tête.
\newtcolorbox{textebox}[1][]{popbase,colframe=popblue,before upper={\popboxhead{Text}{popblue}{#1}\begin{otherlanguage}{english}},after upper={\end{otherlanguage}}}
% Marques ✓ / ✗ (forme correcte / fautive) souvent utilisées par les modèles
\providecommand{\cross}{\textcolor{poppink}{\ensuremath{\times}}}
\providecommand{\xmark}{\cross}\providecommand{\crossmark}{\cross}\providecommand{\cmark}{\ensuremath{\checkmark}}
% Ligne de réponse / justification à compléter (inventée par les modèles)
\providecommand{\justification}[1]{\par\noindent\textit{Justification~:} \dotfill\par}
% \textipa (package tipa non chargé) : la police ne couvre pas l'API -> simple passage
\providecommand{\textipa}[1]{#1}
% Soulignement (ulem non chargé) : \uline, \ul
\providecommand{\uline}[1]{\underline{#1}}\providecommand{\ul}[1]{\underline{#1}}
% \glossaire{\item ...} inventée par les modèles : liste de glossaire
\providecommand{\glossaire}[1]{\begin{itemize}#1\end{itemize}}
% \alert (beamer) inventée par les modèles : texte mis en évidence
\providecommand{\alert}[1]{\textbf{\textcolor{poppink}{#1}}}
% variantes ASCII d'ordinaux écrites en macro
\providecommand{\iere}{\textsuperscript{re}}\providecommand{\ieme}{\textsuperscript{e}}\providecommand{\ere}{\textsuperscript{re}}\providecommand{\eme}{\textsuperscript{e}}\providecommand{\er}{\textsuperscript{er}}\providecommand{\ie}{\textsuperscript{e}}
% Ordinaux français écrits en macro par les modèles : 1\ière, 3\ième
\newcommand{\ière}{\textsuperscript{re}}\newcommand{\ième}{\textsuperscript{e}}
% \exemple (inventée par les modèles) : étiquette « Exemple : »
\providecommand{\exemple}{\textbf{Exemple~:}\ }
% \square utilisable aussi hors mode math (cases à cocher)
\AtBeginDocument{\let\squareMath\square\renewcommand{\square}{\ensuremath{\squareMath}}}
% Mot ou phrase en anglais à l'intérieur d'un paragraphe français
\newcommand{\eng}[1]{\ifmmode\text{\textit{#1}}\else\begingroup\selectlanguage{english}\itshape #1\endgroup\fi}
\let\esp\eng % alias toléré
% flèches d'intonation inventées par les modèles
\providecommand{\textsearrow}{}\renewcommand{\textsearrow}{\ensuremath{\searrow}}\providecommand{\textnearrow}{}\renewcommand{\textnearrow}{\ensuremath{\nearrow}}
% \fr{..} : mot français (inventé par les modèles)
\providecommand{\fr}[1]{\textit{#1}}
\newtcolorbox{exemplebox}[1][]{popbase,colframe=poporange,before upper={\popboxhead{Exemple}{poporange}{#1}}}
\newtcolorbox{definition}[1][]{popbase,colframe=popblue,before upper={\popboxhead{Définition}{popblue}{#1}}}
\newtcolorbox{propriete}[1][]{popbase,colframe=poppurple,before upper={\popboxhead{Propriété}{poppurple}{#1}}}
\newtcolorbox{methode}[1][]{popbase,colframe=popgold,before upper={\popboxhead{Méthode}{popgold}{#1}}}
\newtcolorbox{attention}[1][]{popbase,colframe=poppink,before upper={\popboxhead{Attention}{poppink}{#1}}}
\newtcolorbox{experience}[1][]{popbase,colframe=popblue,colback=popblueL,before upper={\popboxhead{Expérience}{popblue}{#1}}}
% « Le savais-tu ? » : carte violette pleine (contexte, culture, Cameroun)
\newtcolorbox{savaistu}[1][]{popbase,colback=poppurple,colframe=popink,
  fontupper=\popsemi\fontsize{10.4}{14.2}\selectfont\color{white},
  before upper={\poppill{Le savais-tu\,?}{white}{poppurple}\ifx\relax#1\relax\else\hspace{6pt}{\popxbold\color{white}#1}\fi\par\vspace{7pt}}}
% Exercice résolu : énoncé en haut, \tcblower puis la solution sur fond crème
\newcounter{popexo}\newcounter{popexr}
\newtcolorbox{exoresolu}[1][]{popbase,colframe=poporange,colbacklower=popcream,
  segmentation style={popink,line width=1.2pt,dashed},
  before upper={\stepcounter{popexr}\popboxhead{Exercice résolu \thepopexr}{poporange}{#1}},
  before lower={\poppill{Solution}{popink}{popcream}\par\vspace{6pt}}}
% Exercice (non résolu en place) — la correction est dans la partie Corrigés
\newtcolorbox{exercice}[1][]{popbase,colframe=popink,
  before upper={\stepcounter{popexo}\popboxhead{Exercice \thepopexo}{popink}{#1}}}
% Corrigé
\newtcolorbox{corrige}[1][]{popbase,colframe=popgreen,colback=popgreenL,
  before upper={\popboxhead{Corrigé}{popgreen}{#1}}}
% Objectifs / prérequis / bilan
\newtcolorbox{objectifs}{popbase,colframe=popink,colback=poporangeL,
  before upper={\popboxhead{Objectifs de la leçon}{popink}{}}}
\newtcolorbox{prerequis}{popbase,colframe=popink,colback=popblueL,
  before upper={\popboxhead{Prérequis}{popblue}{}}}
\newtcolorbox{bilan}{popbase,colframe=popink,colback=white,boxrule=2.8pt,
  before upper={\popboxhead{Fiche bilan}{poporange}{L'essentiel à connaître pour l'examen}}}
% Figure encadrée avec légende
\newtcolorbox{popfigure}[1][]{enhanced,breakable=false,colback=white,colframe=popline,boxrule=1.4pt,arc=8pt,
  left=10pt,right=10pt,top=10pt,bottom=8pt,before skip=10pt,after skip=12pt,halign=center,
  lower separated=false,halign lower=center,
  fontlower=\popsemi\fontsize{9}{11.5}\selectfont\color{popmuted},
  #1}
\newcommand{\legende}[1]{\par\vspace{6pt}{\popsemi\fontsize{9}{11.5}\selectfont\color{popmuted}#1}}

% ============================================================
% Signature OnBuch+
% ============================================================
\newcommand{\poplogoinline}[1]{{\poplogo\fontsize{#1}{#1}\selectfont\color{popink}OnBuch\color{poporange}+}}
\newcommand{\signatureonbuch}{%
  \begin{tcolorbox}[enhanced,colback=popink,colframe=popink,boxrule=2.2pt,arc=10pt,
      drop shadow=poporange,left=14pt,right=14pt,top=10pt,bottom=10pt,before skip=0pt,after skip=0pt]
    \tikz[baseline=-0.7ex]\node[circle,fill=popgreen,draw=popcream,line width=1.3pt,minimum size=22pt,inner sep=0]{\popcheck{white}};%
    \hspace{9pt}\begin{minipage}[c]{0.62\linewidth}
      {\popxbold\fontsize{12}{13}\selectfont\color{popcream}Signé\ \textbullet\ {\poplogo OnBuch\color{poporange}+}}\par\vspace{2pt}
      {\raggedright\popsemi\fontsize{8}{10.5}\selectfont\color{popline}\MakeUppercase{Équipe pédagogique OnBuch+}\\\MakeUppercase{Conforme au programme officiel MINESEC}\par}
    \end{minipage}\hfill
    {\poplogo\fontsize{22}{22}\selectfont\color{popcream}OnBuch\color{poporange}+}
  \end{tcolorbox}%
}

% ============================================================
% Page de garde
% #1 titre · #2 matière · #3 niveau · #4 module · #5 n° leçon · #6 durée ·
% #7 description / objectifs (texte ou itemize)
% ============================================================
\newcommand{\pagegarde}[7]{%
  \thispagestyle{empty}
  \newgeometry{a4paper,margin=1.7cm,top=1.6cm,bottom=1.4cm}%
  \begin{tikzpicture}[remember picture,overlay]
    \fill[poporange!18] ([xshift=-3.2cm,yshift=-3.6cm]current page.north east) circle (4.6cm);
    \fill[poppurple!14] ([xshift=2.4cm,yshift=7.6cm]current page.south west) circle (3.8cm);
  \end{tikzpicture}%
  {\poplogo\fontsize{30}{30}\selectfont\color{popink}OnBuch\color{poporange}+}\hfill
  \raisebox{0.6ex}{\poppill{Cours signé \textbullet\ Premium}{popink}{popcream}}\par
  \vspace{1pt}\popsquiggle{poppurple}\par
  \vspace{8pt}
  \begin{tcolorbox}[enhanced,colback=poporange,colframe=popink,boxrule=2.8pt,arc=13pt,
      drop shadow=popink,left=18pt,right=18pt,top=12pt,bottom=13pt,after skip=10pt]
    \poppill{#2 \textbullet\ #3}{popink}{popcream}\hspace{5pt}\poppill{#4}{white}{popink}\par\vspace{8pt}
    {\popdisplay\fontsize{14}{14}\selectfont\color{popink}LEÇON #5}\par\vspace{4pt}
    {\raggedright\hyphenpenalty=10000\exhyphenpenalty=10000\popdisplay\fontsize{25}{27}\selectfont\color{popcream}\MakeUppercase{#1}\par}
    \vspace{8pt}
    {\popxbold\fontsize{9.5}{9.5}\selectfont\color{popink}\MakeUppercase{Durée conseillée : #6}}
  \end{tcolorbox}
  \begin{tcolorbox}[enhanced,colback=white,colframe=popink,boxrule=2.2pt,arc=10pt,
      drop shadow=popink,left=16pt,right=16pt,top=10pt,bottom=10pt,after skip=8pt]
    \poppill{Dans cette leçon}{poppurple}{white}\par\vspace{5pt}
    \popsemi\fontsize{9.3}{12.6}\selectfont\color{popink}#7
  \end{tcolorbox}
  \noindent\begin{minipage}[t]{0.487\linewidth}
    \begin{tcolorbox}[enhanced,colback=popgreen,colframe=popink,boxrule=2.2pt,arc=10pt,
        drop shadow=popink,left=11pt,right=11pt,top=10pt,bottom=10pt,height=3.1cm,valign=center]
      \tcbox[colback=white,colframe=popink,boxrule=1.3pt,arc=5pt,boxsep=0pt,nobeforeafter,
        left=3pt,right=3pt,top=3pt,bottom=3pt,valign=center]{\qrcode[height=1.9cm]{https://play.google.com/store/apps/details?id=com.luvvix.onbuch&hl=fr}}%
      \hspace{8pt}\begin{minipage}[c]{0.5\linewidth}
        {\popdisplay\fontsize{11}{12.5}\selectfont\color{white}\MakeUppercase{Télécharge OnBuch}}\par\vspace{4pt}
        {\raggedright\popsemi\fontsize{8.4}{11}\selectfont\color{white}Gratuit \textbullet\ Google Play\\Tuteur IA, annales, concours, résultats\par}
      \end{minipage}
    \end{tcolorbox}
  \end{minipage}\hfill
  \begin{minipage}[t]{0.487\linewidth}
    \begin{tcolorbox}[enhanced,colback=poppurple,colframe=popink,boxrule=2.2pt,arc=10pt,
        drop shadow=popink,left=11pt,right=11pt,top=10pt,bottom=10pt,height=3.1cm,valign=center]
      \tikz[baseline=-0.7ex]\node[circle,fill=white,draw=popink,line width=1.3pt,minimum size=1.6cm,inner sep=0]{\popdisplay\fontsize{24}{24}\selectfont\color{poppurple}+};%
      \hspace{8pt}\begin{minipage}[c]{0.6\linewidth}
        {\popdisplay\fontsize{11}{12.5}\selectfont\color{white}\MakeUppercase{Passe à OnBuch+}}\par\vspace{4pt}
        {\raggedright\popsemi\fontsize{8.4}{11}\selectfont\color{white}Tous les cours signés, Tuteur IA illimité, fiches PDF — onglet OnBuch+ de l'app\par}
      \end{minipage}
    \end{tcolorbox}
  \end{minipage}
  \vspace{8pt}
  \popcontenu
  \vfill
  \signatureonbuch
  \restoregeometry
  \newpage
}

% Rangée « Ce cours contient » (page de garde)
\newcommand{\poptile}[3]{% #1 couleur #2 gros texte #3 légende
  \begin{tcolorbox}[enhanced,colback=#1,colframe=popink,boxrule=2pt,arc=9pt,shadow={2.5pt}{-2.5pt}{0pt}{popink},
      left=6pt,right=6pt,top=9pt,bottom=9pt,halign=center,valign=center,height=1.75cm,width=\linewidth,nobeforeafter]
    {\popdisplay\fontsize{12.5}{13}\selectfont\color{white}\MakeUppercase{#2}}\par\vspace{3pt}
    {\centering\popsemi\fontsize{7.8}{9.5}\selectfont\color{white}#3\par}
  \end{tcolorbox}}
\newcommand{\popcontenu}{%
  \par\noindent{\popxboldls\fontsize{7.5}{7.5}\selectfont\color{popmuted}CE COURS SIGNÉ CONTIENT}\par\vspace{7pt}
  \noindent\begin{minipage}[t]{0.235\linewidth}\poptile{popblue}{Cours}{complet, illustré, pas à pas}\end{minipage}\hfill
  \begin{minipage}[t]{0.235\linewidth}\poptile{poporange}{Méthodes}{et exercices résolus}\end{minipage}\hfill
  \begin{minipage}[t]{0.235\linewidth}\poptile{poppink}{Exercices}{type examen + corrigés}\end{minipage}\hfill
  \begin{minipage}[t]{0.235\linewidth}\poptile{popgreen}{Fiche bilan}{l'essentiel à retenir}\end{minipage}\par}

% Sommaire
\newtcolorbox{sommaire}{enhanced,colback=white,colframe=popink,boxrule=2.2pt,arc=10pt,
  drop shadow=popink,left=18pt,right=18pt,top=13pt,bottom=13pt,before skip=6pt,after skip=20pt,
  before upper={\poppill{Sommaire}{poppurple}{white}\par\vspace{9pt}\popsemi\fontsize{10.6}{14.5}\selectfont\color{popink}}}

% Signature de fin de cours — poussée tout en bas de la dernière page
\newcommand{\findecours}{%
  \par\vfill\nopagebreak
  \begin{center}\popsquiggle{poporange}\end{center}\vspace{4pt}
  \signatureonbuch
}
\AtBeginDocument{\def\llbracket{\mathopen{[\mkern-3.5mu[}}\def\rrbracket{\mathclose{]\mkern-3.5mu]}}} % intervalles d'entiers (glyphe ⟦ absent de la police)
% Glyphes absents de Fira Math : redéfinis à partir de glyphes disponibles
\AtBeginDocument{%
  \def\setminus{\mathbin{\backslash}}\let\smallsetminus\setminus
  \def\vdots{\mathord{\vbox{\baselineskip3.2pt\lineskiplimit0pt\kern4pt\hbox{.}\hbox{.}\hbox{.}}}}%
  \def\ddots{\mathinner{\mkern1mu\raise6.5pt\hbox{.}\mkern2mu\raise3.6pt\hbox{.}\mkern2mu\raise0.7pt\hbox{.}\mkern1mu}}%
  \def\triangle{\mathord{\scalebox{0.9}{$\symup{\Delta}$}}}%
  \def\nabla{\mathord{\raisebox{\depth}{\scalebox{1}[-1]{$\symup{\Delta}$}}}}%
  \def\checkmark{\popcheck{popdark}}%
  \def\star{\mathbin{\ast}}\def\bigstar{\mathord{\ast}}%
  \def\diamond{\mathbin{\scalebox{0.6}{$\Diamond$}}}%
  \def\bigcup{\mathop{\mathchoice{\raisebox{-0.3em}{\scalebox{1.7}{$\cup$}}}{\scalebox{1.25}{$\cup$}}{\cup}{\cup}}\displaylimits}%
  \def\bigcap{\mathop{\mathchoice{\raisebox{-0.3em}{\scalebox{1.7}{$\cap$}}}{\scalebox{1.25}{$\cap$}}{\cap}{\cap}}\displaylimits}%
  \def\lesseqqgtr{\mathrel{\lessgtr}}\def\bigoplus{\mathop{\oplus}\displaylimits}%
}
\providecommand{\ssi}{si et seulement si} % abréviation fréquente
\AtBeginDocument{\providecommand{\wideparen}{\overparen}} % arc de cercle

\begin{document}

\pagegarde{Lesson 4 — Speaking: Exchanges information about applying for a passport}{Anglais}{Tle A}{Chapter 1 : Family and Social Life}{EN04}{2 h}{Cette leçon d'expression orale apprend à l'élève de Terminale A à échanger des informations en anglais dans le cadre d'une demande de passeport : poser des questions, donner des renseignements personnels, relancer et conclure poliment un dialogue administratif. Elle fournit des dialogues modèles, des formules fonctionnelles et des conseils de prononciation pour des jeux de rôle guidés.
\vspace{4pt}
\begin{multicols}{2}\small
\begin{itemize}
\item À la fin de cette leçon, je sais saluer et me présenter poliment à un guichet administratif en anglais.
\item À la fin de cette leçon, je sais demander des informations sur les documents, les frais et les délais d'obtention d'un passeport.
\item À la fin de cette leçon, je sais donner mes informations personnelles (nom, date et lieu de naissance, adresse, profession) en anglais correct.
\item À la fin de cette leçon, je sais utiliser les formules fonctionnelles pour demander, remercier, relancer et conclure un échange.
\item À la fin de cette leçon, je sais faire répéter ou préciser une information que je n'ai pas comprise.
\item À la fin de cette leçon, je sais prononcer correctement le vocabulaire administratif clé avec la bonne accentuation et la bonne intonation.
\end{itemize}\end{multicols}}

\begin{sommaire}
\begin{multicols}{2}
\begin{enumerate}
\item Découverte : la situation et le vocabulaire clé
\item Dialogue modèle 1 : demander des informations
\item Dialogue modèle 2 : déposer sa demande
\item Boîte à outils : les formules fonctionnelles
\item Prononciation, accentuation et intonation
\item Jeux de rôle guidés et production libre
\item Activité d'intégration
\item Méthodes et exercices résolus
\item Exercices
\item Corrigés des exercices
\item Fiche bilan
\end{enumerate}
\end{multicols}
\end{sommaire}

\begin{objectifs}
\begin{itemize}
\item À la fin de cette leçon, je sais saluer et me présenter poliment à un guichet administratif en anglais.
\item À la fin de cette leçon, je sais demander des informations sur les documents, les frais et les délais d'obtention d'un passeport.
\item À la fin de cette leçon, je sais donner mes informations personnelles (nom, date et lieu de naissance, adresse, profession) en anglais correct.
\item À la fin de cette leçon, je sais utiliser les formules fonctionnelles pour demander, remercier, relancer et conclure un échange.
\item À la fin de cette leçon, je sais faire répéter ou préciser une information que je n'ai pas comprise.
\item À la fin de cette leçon, je sais prononcer correctement le vocabulaire administratif clé avec la bonne accentuation et la bonne intonation.
\item À la fin de cette leçon, je sais jouer un dialogue complet « demandeur / agent » avec un camarade, sans lire mes notes.
\item À la fin de cette leçon, je sais adapter mon niveau de politesse (Could you...?, I'd like...) à une situation formelle.
\end{itemize}
\end{objectifs}

\begin{prerequis}
\begin{itemize}
\item Les questions en Wh- (What, Where, When, How much, How long) et en Yes/No (Do you...? / Can I...?) vues en Première.
\item Le vocabulaire de l'identité et de la famille (name, surname, date of birth, marital status) vu en début d'année dans le chapitre Family and Social Life.
\item Les formules de politesse de base (please, thank you, excuse me) et les salutations formelles.
\item Les nombres, les dates et l'heure en anglais (pour les frais, délais et rendez-vous).
\item Le présent simple et « would like » pour exprimer une demande polie.
\end{itemize}
\end{prerequis}

\newpage

\coursec{Découverte : la situation et le vocabulaire clé}

\courssub{La situation d'accroche}

Imagine : tu viens d'être admis dans une université au Ghana ou tu veux rendre visite à ta famille à Bamenda, et on te demande un passeport. À la Délégation Générale de la Sûreté Nationale à Yaoundé, ou dans une ambassade, tout se passe souvent en anglais : formulaires, questions de l'agent, rendez-vous. Sauras-tu demander les documents nécessaires, donner tes informations personnelles et comprendre les réponses ? Cette leçon t'apprend à mener ce dialogue de A à Z, comme un vrai candidat au Bac.

\textbf{Problématique.} Comment échanger des informations en anglais pour faire une demande de passeport : quels mots employer, quelles questions poser, quelles réponses comprendre ?

\courssub{Observation : un texte pour découvrir la situation}

Avant d'apprendre les règles, lis ce court texte. Il décrit, en anglais simple, le parcours d'une élève de terminale qui demande son premier passeport. Souligne mentalement tous les mots qui désignent des \cle{documents}, des \cle{personnes} et des \cle{actions}.

\begin{textebox}[At the passport office — a reading text]
Last month, Divine, a student in Terminale at a high school in Yaoundé, decided to apply for a passport. She wants to study nursing in Ghana next year, and the university asked her for a valid passport before the end of the school year.

First, Divine went to the police station near her home to ask for information. A clerk at the counter explained the procedure to her in English. She needed four documents: a birth certificate, a national ID card, two passport-sized photographs and a receipt of payment. The fee was not very high, but she had to pay it at the bank before submitting her file.

The following week, Divine filled in the application form carefully. She wrote her full name, her date and place of birth, her address and her parents' names. Then she submitted the form and the documents at the counter of the immigration office. The officer checked everything, took her fingerprints and gave her an appointment slip.

Two weeks later, Divine received a phone call: her passport was ready. She went back to the office and collected it at the same counter. The immigration officer smiled and said: “Your passport is valid for five years. Do not forget to renew it before it expires.” Divine was very proud. Her dream of studying abroad was becoming real, step by step.
\end{textebox}

\textbf{Glossaire.}

\begin{center}
\begin{tabularx}{\linewidth}{|l|l|X|}\hline
\rowcolor{popblueL} \textbf{Mot anglais} & \textbf{Nature} & \textbf{Traduction} \\ \hline
to apply for & verbe & faire une demande de \\ \hline
valid & adjectif & valide, en cours de validité \\ \hline
a clerk & nom & un employé de bureau, un guichetier \\ \hline
a receipt of payment & nom & un reçu de paiement \\ \hline
a fee & nom & des frais (officiels) \\ \hline
to fill in & verbe & remplir (un formulaire) \\ \hline
to submit & verbe & soumettre, déposer \\ \hline
fingerprints & nom & empreintes digitales \\ \hline
an appointment slip & nom & une convocation, un bordereau de rendez-vous \\ \hline
to collect & verbe & retirer, venir chercher \\ \hline
to renew & verbe & renouveler \\ \hline
to expire & verbe & expirer, arriver à échéance \\ \hline
\end{tabularx}
\end{center}

\textbf{Questions de compréhension (oral).}
\begin{enumerate}
\item Why does Divine need a passport?
\item What four documents does she need?
\item Where does she submit her application form?
\item How long does she wait before collecting her passport?
\item What must she do before the passport expires?
\end{enumerate}

\begin{corrige}[Questions de compréhension]
\begin{enumerate}
\item \eng{She needs a passport because she wants to study nursing in Ghana, and the university asked her for a valid passport.}
\item \eng{She needs a birth certificate, a national ID card, two passport-sized photographs and a receipt of payment.}
\item \eng{She submits it at the counter of the immigration office.}
\item \eng{She waits two weeks.}
\item \eng{She must renew it before it expires.}
\end{enumerate}
\end{corrige}

\courssub{Brainstorming : quels documents faut-il ?}

Observe le texte : Divine prépare quatre documents. En classe, le professeur peut écrire au tableau la question \eng{What documents do you need to apply for a passport?} et noter toutes les réponses des élèves. Voici les réponses attendues :

\begin{itemize}
\item \eng{a birth certificate} — un acte de naissance ;
\item \eng{a national ID card} — une carte nationale d'identité ;
\item \eng{two passport-sized photographs} — deux photos d'identité (format passeport) ;
\item \eng{a receipt of payment} — un reçu de paiement (la preuve que les frais ont été payés).
\end{itemize}

\begin{textebox}[Key sentence — la phrase clé du texte]
To apply for a passport in Cameroon, you need a birth certificate, a national ID card, two passport-sized photographs and a receipt of payment. You fill in an application form, submit it at the counter, and collect your passport about two weeks later.
\end{textebox}

Traduction : « Pour faire une demande de passeport au Cameroun, il faut un acte de naissance, une carte nationale d'identité, deux photos d'identité et un reçu de paiement. Tu remplis un formulaire de demande, tu le déposes au guichet, et tu retires ton passeport environ deux semaines plus tard. »

\courssub{Champ lexical 1 : les documents}

Voici le vocabulaire essentiel des documents. Apprends chaque mot avec son article et sa prononciation approximative en lettres françaises.

\begin{center}
\begin{tabularx}{\linewidth}{|l|l|X|}\hline
\rowcolor{popblueL} \textbf{English} & \textbf{Prononciation} & \textbf{Français} \\ \hline
a passport & « è pa-ce-porte » & un passeport \\ \hline
an application form & « è-n a-pli-ké-cheune forme » & un formulaire de demande \\ \hline
a birth certificate & « e beurthe seur-ti-fi-kète » & un acte de naissance \\ \hline
an ID card (identity card) & « è-n aï-di karde » & une carte d'identité \\ \hline
a passport-sized photograph & « è pa-ce-porte-saïzd fo-to-grafe » & une photo d'identité format passeport \\ \hline
a receipt & « è ri-site » (le p ne se prononce pas) & un reçu \\ \hline
a visa & « è vi-za » & un visa \\ \hline
\end{tabularx}
\end{center}

\begin{attention}[Orthographe et prononciation : receipt]
Le mot \eng{receipt} s'écrit avec un \textbf{p} qui ne se prononce pas : « ri-site ». De même, on ne prononce pas le \textbf{b} de \eng{debt} (dette). Ne confonds pas \eng{receipt} (le reçu) et \eng{recipe} (la recette de cuisine) : « ré-ssi-pi ».
\end{attention}

\courssub{Champ lexical 2 : les personnes et les lieux}

\begin{center}
\begin{tabularx}{\linewidth}{|l|X|}\hline
\rowcolor{popblueL} \textbf{English} & \textbf{Français} \\ \hline
an applicant & un demandeur, un candidat (celui qui fait la demande) \\ \hline
an immigration officer & un agent de l'immigration \\ \hline
a clerk & un employé de bureau, un guichetier \\ \hline
an embassy & une ambassade \\ \hline
a consulate & un consulat \\ \hline
a police station & un commissariat de police \\ \hline
a counter & un guichet \\ \hline
\end{tabularx}
\end{center}

\begin{definition}[An applicant]
\eng{An applicant} = a person who officially asks for something (un demandeur, un candidat). Le verbe correspondant est \eng{to apply} et le nom de l'action est \eng{an application}.\par
\eng{The applicant must sign the form.} = « Le demandeur doit signer le formulaire. »
\end{definition}

\begin{attention}[Faux ami : « a demand » n'est pas « une demande »]
En anglais, \eng{a demand} signifie « une exigence, une revendication », jamais « une demande administrative ». Pour dire « faire une demande de passeport », on dit \eng{to apply for a passport} ou \eng{to make an application for a passport}.\par
\eng{She made an application for a visa.} = « Elle a fait une demande de visa. »\par
\eng{The workers' demands were accepted.} = « Les revendications des travailleurs ont été acceptées. »
\end{attention}

\courssub{Champ lexical 3 : les actions}

Les verbes sont le moteur du dialogue. Observe-les d'abord dans des phrases tirées du texte, puis apprends le tableau.

\begin{center}
\begin{tabularx}{\linewidth}{|l|l|X|}\hline
\rowcolor{popblueL} \textbf{Verbe anglais} & \textbf{Français} & \textbf{Exemple} \\ \hline
to apply for & demander (officiellement) & \eng{She applied for a passport.} \\ \hline
to fill in a form & remplir un formulaire & \eng{Fill in this form, please.} \\ \hline
to submit & soumettre, déposer & \eng{He submitted his file on Monday.} \\ \hline
to pay a fee & payer des frais & \eng{You must pay the fee at the bank.} \\ \hline
to collect & retirer, venir chercher & \eng{Come back and collect your passport.} \\ \hline
to renew & renouveler & \eng{I must renew my ID card.} \\ \hline
to expire & expirer & \eng{My passport expires next year.} \\ \hline
\end{tabularx}
\end{center}

\begin{propriete}[Trois verbes à particule à ne pas confondre]
L'anglais utilise des verbes suivis d'une petite particule (\eng{for}, \eng{in}, \eng{up}) qui change le sens. Trois distinctions sont vitales dans notre situation :
\begin{itemize}
\item \eng{to apply \textbf{for} a passport} = demander un passeport (\eng{for} introduit la chose demandée) ;
\item \eng{to fill \textbf{in} a form} = remplir un formulaire (en anglais britannique, on dit aussi \eng{to fill \textbf{out} a form}, surtout en anglais américain) ;
\item \eng{to pick \textbf{up}} ou \eng{to collect} = retirer, venir chercher le document une fois prêt.
\end{itemize}
\eng{I applied for a passport, I filled in the form, and two weeks later I picked it up.} = « J'ai demandé un passeport, j'ai rempli le formulaire, et deux semaines plus tard je l'ai retiré. »
\end{propriete}

\begin{exemplebox}[Exemples traduits à réutiliser à l'oral]
\eng{I would like to apply for a passport.} = « Je voudrais faire une demande de passeport. »\par
\eng{How much is the fee?} = « Quel est le montant des frais ? »\par
\eng{When can I collect it?} = « Quand puis-je le retirer ? »\par
\eng{Where do I submit the form?} = « Où est-ce que je dépose le formulaire ? »\par
\eng{My passport has expired; I need to renew it.} = « Mon passeport a expiré ; je dois le renouveler. »
\end{exemplebox}

\courssub{Carte mentale et frise des étapes}

\begin{popfigure}
\begin{tikzpicture}[mindmap, grow cyclic, every node/.style={concept}, text=white,
root concept/.append style={concept color=popblue, font=\bfseries},
level 1/.append style={level distance=3.2cm, sibling angle=90}]
\node[root concept]{APPLYING FOR A PASSPORT}
child[concept color=poporange]{ node {Documents}
child[grow=180, concept color=poporange!70]{ node[font=\small, align=center]{birth certificate\\ID card\\photos\\receipt} }}
child[concept color=popgreen]{ node {People \& Places}
child[grow=90, concept color=popgreen!70]{ node[font=\small, align=center]{applicant\\officer\\embassy\\counter} }}
child[concept color=poppurple]{ node {Actions}
child[grow=0, concept color=poppurple!70]{ node[font=\small, align=center]{apply for\\fill in\\submit\\collect} }}
child[concept color=popgold]{ node {Money \& Time}
child[grow=270, concept color=popgold!70]{ node[font=\small, align=center]{pay a fee\\two weeks\\valid\\expires} }};
\end{tikzpicture}
\legende{Carte mentale du vocabulaire : les quatre familles de mots de la demande de passeport.}
\end{popfigure}

\begin{popfigure}
\begin{tikzpicture}[node distance=0.35cm,
etape/.style={rectangle, rounded corners, draw=popblue, fill=popblueL, text width=2.5cm, align=center, font=\small, minimum height=1cm},
fleche/.style={-{Stealth[length=2mm]}, thick, popdark}]
\node[etape] (e1) {Ask for information};
\node[etape, right=of e1] (e2) {Fill in the form};
\node[etape, right=of e2] (e3) {Submit documents};
\node[etape, right=of e3] (e4) {Pay the fee};
\node[etape, right=of e4] (e5) {Wait};
\node[etape, right=of e5] (e6) {Collect the passport};
\draw[fleche] (e1) -- (e2);
\draw[fleche] (e2) -- (e3);
\draw[fleche] (e3) -- (e4);
\draw[fleche] (e4) -- (e5);
\draw[fleche] (e5) -- (e6);
\end{tikzpicture}
\legende{Les six étapes de la demande de passeport : apprends-les dans l'ordre, elles structurent le dialogue.}
\end{popfigure}

\courssub{Comparaison français / anglais : les pièges du lexique}

\begin{center}
\begin{tabularx}{\linewidth}{|X|X|X|}\hline
\rowcolor{popblueL} \textbf{Français} & \textbf{English} & \textbf{Remarque} \\ \hline
faire une demande & to apply / to make an application & jamais \eng{to make a demand} (exiger) \\ \hline
un formulaire & a form & \eng{formula} = une formule (maths, chimie) \\ \hline
remplir un formulaire & to fill in a form & la particule \eng{in} est obligatoire en anglais britannique \\ \hline
retirer son passeport & to collect / to pick up one's passport & \eng{to remove} = enlever, ôter \\ \hline
des frais & a fee / fees & \eng{fresh} = frais (adjectif) ! \\ \hline
un acte de naissance & a birth certificate & mot à mot : « certificat de naissance » \\ \hline
valide / expiré & valid / expired & \eng{My passport has expired.} \\ \hline
\end{tabularx}
\end{center}

\courssub{Exercice résolu et entraînement}

\begin{exoresolu}[Complète le dialogue au guichet]
Énoncé. Complète ce mini-dialogue avec les mots suivants : \eng{apply for — fill in — receipt — collect — fee}.\par
\eng{Clerk: Good morning. Can I help you?}\par
\eng{Applicant: Yes, I would like to \ldots\ldots\ldots\ldots a passport.}\par
\eng{Clerk: Please \ldots\ldots\ldots\ldots this application form and bring your birth certificate, your ID card and the \ldots\ldots\ldots\ldots of payment.}\par
\eng{Applicant: How much is the \ldots\ldots\ldots\ldots ?}\par
\eng{Clerk: It is written on the notice board. You can \ldots\ldots\ldots\ldots your passport two weeks after you submit the file.}
\tcblower
Solution détaillée.
\begin{enumerate}
\item \eng{apply for} : après \eng{I would like to}, on met la base verbale ; on « demande » un passeport, donc \eng{apply for}.
\item \eng{fill in} : le guichetier donne un ordre (impératif) ; on « remplit » un formulaire avec \eng{fill in}.
\item \eng{receipt} : la preuve du paiement est \eng{the receipt of payment} (attention, « ri-site », le p est muet).
\item \eng{fee} : la question \eng{How much...?} porte sur le prix, donc \eng{the fee} (les frais).
\item \eng{collect} : après deux semaines d'attente, on « retire » le passeport : \eng{collect}.
\end{enumerate}
Dialogue complet : \eng{I would like to apply for a passport. — Please fill in this application form and bring your birth certificate, your ID card and the receipt of payment. — How much is the fee? — You can collect your passport two weeks after you submit the file.}
\end{exoresolu}

\begin{exercice}[À toi de jouer]
\begin{enumerate}
\item Traduis en anglais : « Je voudrais renouveler mon passeport, car il a expiré. »
\item Trouve l'erreur et corrige : \eng{I made a demand for a visa at the ambassy.}
\item Classe ces mots dans les quatre colonnes de la carte mentale : \eng{clerk — visa — to renew — consulate — to expire — birth certificate — fee — to submit}.
\item Réponds à la question du guichetier : \eng{What documents do you need?} (réponds avec une phrase complète commençant par \eng{I need...}).
\end{enumerate}
\end{exercice}

\begin{corrige}[Exercice « À toi de jouer »]
\begin{enumerate}
\item \eng{I would like to renew my passport because it has expired.} (On peut aussi dire : \eng{...because it expired last month.})
\item Deux erreurs : \eng{I made an application for a visa at the embassy.} — \eng{a demand} est un faux ami (« une exigence ») et \eng{embassy} prend un \textbf{e}, pas un \textbf{a}.
\item Documents : \eng{visa, birth certificate} — People \& Places : \eng{clerk, consulate} — Actions : \eng{to renew, to submit} — Money \& Time : \eng{fee, to expire}.
\item \eng{I need a birth certificate, a national ID card, two passport-sized photographs and a receipt of payment.}
\end{enumerate}
\end{corrige}

\begin{savaistu}[Le passeport au Cameroun et dans le monde anglophone]
Au Cameroun, le passeport biométrique se demande d'abord en ligne (pré-enrôlement), puis le demandeur se présente physiquement à la Délégation Générale de la Sûreté Nationale (DGSN) à Yaoundé ou dans un commissariat central pour l'enrôlement, avant de retirer le document. Au Royaume-Uni, l'administration qui délivre les passeports s'appelle le \eng{passport office} ; aux États-Unis, on parle de \eng{passport agency}. À l'étranger, un Camerounais qui perd son passeport s'adresse à l'ambassade (\eng{embassy}) ou au consulat (\eng{consulate}) de son pays. Dans tous ces lieux, la première phrase à connaître est : \eng{I would like to apply for a passport.}
\end{savaistu}

\begin{aretenir}[L'essentiel de la section]
\begin{itemize}
\item La situation : demander un passeport (\eng{to apply for a passport}) à la DGSN, au commissariat ou dans une ambassade.
\item Quatre documents clés : \eng{birth certificate, national ID card, passport-sized photographs, receipt of payment}.
\item Les personnes et lieux : \eng{applicant, clerk, immigration officer, embassy, consulate, police station, counter}.
\item Les actions : \eng{to apply for, to fill in, to submit, to pay a fee, to collect, to renew, to expire}.
\item Piège majeur : \eng{a demand} = une exigence ; « une demande » se dit \eng{an application}.
\item Les six étapes dans l'ordre : \eng{ask for information → fill in the form → submit documents → pay the fee → wait → collect the passport}.
\end{itemize}
\end{aretenir}

\coursec{Dialogue modèle 1 : demander des informations}

Dans la section précédente, nous avons découvert la situation de communication — un citoyen qui se rend au poste de police ou au centre de délivrance des passeports pour obtenir des informations — ainsi que le vocabulaire clé : \eng{passport}, \eng{birth certificate}, \eng{ID card}, \eng{passport-sized photographs}, \eng{receipt of payment}, \eng{fee}, \eng{to apply for}. Nous allons maintenant écouter, lire et analyser un premier dialogue modèle entre un \cle{demandeur} (\eng{Applicant}) et un \cle{agent} (\eng{Officer}) au guichet. Ce dialogue sera notre référence : c'est sur lui que nous repérerons les étapes d'un échange administratif et les formules fonctionnelles à réutiliser ensuite dans vos propres jeux de rôle.

\courssub{Écoute et lecture du dialogue modèle}

Fermez d'abord vos cahiers : écoutez le dialogue lu par le professeur (ou joué par deux élèves) et essayez de répondre mentalement à trois questions : \emph{Where are the speakers? What does the applicant want? How much must he pay?} Puis ouvrez vos cahiers et lisez le dialogue en suivant les formules surlignées.

\begin{textebox}[At the passport office — Dialogue 1]
\textbf{Officer:} Good morning, can I help you?

\textbf{Applicant:} Good morning. I'd like to apply for a passport, please.

\textbf{Officer:} Certainly. Do you have your birth certificate and your ID card?

\textbf{Applicant:} Yes, here they are. What else do I need?

\textbf{Officer:} Two passport-sized photographs and the receipt of payment.

\textbf{Applicant:} How much is the fee?

\textbf{Officer:} It's 75,000 francs.

\textbf{Applicant:} And how long does it take?

\textbf{Officer:} About two weeks.

\textbf{Applicant:} Thank you very much.

\textbf{Officer:} You're welcome.
\end{textebox}

\textbf{Glossaire}

\begin{center}
\begin{tabularx}{\linewidth}{|l|X|l|}\hline
\rowcolor{popblueL} \textbf{English} & \textbf{Nature} & \textbf{Français} \\ \hline
to apply for (a passport) & verbe & faire une demande de (passeport) \\ \hline
birth certificate & nom & acte de naissance \\ \hline
ID card (identity card) & nom & carte nationale d'identité \\ \hline
passport-sized photograph & nom & photo d'identité (format passeport) \\ \hline
receipt of payment & nom & reçu de paiement, quittance \\ \hline
fee & nom & frais (administratifs) \\ \hline
How long does it take? & expression & Combien de temps cela prend-il ? \\ \hline
\end{tabularx}
\end{center}

\begin{aretenir}[Prononciation des mots clés]
\begin{itemize}
\item \eng{passport} se prononce « \textbf{passe}-porte » (l'accent est sur la première syllabe ; le second \emph{a} est très faible, \eng{/ˈpɑːspɔːt/}).
\item \eng{certificate} se prononce « \textbf{ser}-ti-fi-ket » (quatre syllabes, accent sur « ti », \eng{/səˈtɪfɪkət/}).
\item \eng{receipt} se prononce « ri-\textbf{ssite} » (\eng{/rɪˈsiːt/}) : attention, le \emph{p} ne se prononce pas !
\item \eng{photograph} se prononce « \textbf{fo}-to-graf » (accent sur la première syllabe, \eng{/ˈfəʊtəɡrɑːf/}).
\item \eng{fee} se prononce « \textbf{fii} » (une seule syllabe longue, \eng{/fiː/}).
\end{itemize}
\end{aretenir}

\courssub{Repérage des étapes du dialogue}

Tout dialogue administratif suit une progression logique et prévisible. Repérer ces étapes, c'est déjà comprendre la moitié du dialogue. Observons le nôtre :

\begin{enumerate}
\item \textbf{Salutation} (\eng{Greeting}) : \eng{Good morning, can I help you?} — l'agent ouvre l'échange et propose son aide.
\item \textbf{Annonce de l'objet de la visite} (\eng{Purpose}) : \eng{I'd like to apply for a passport, please.} — le demandeur dit clairement ce qu'il veut.
\item \textbf{Questions sur les documents} (\eng{Questions}) : \eng{Do you have your birth certificate...? What else do I need?}
\item \textbf{Réponses et informations chiffrées} (\eng{Answers}) : les documents requis, le montant des frais (\eng{It's 75,000 francs}), le délai (\eng{About two weeks}).
\item \textbf{Remerciements et conclusion} (\eng{Thanks \& Goodbye}) : \eng{Thank you very much. — You're welcome.}
\end{enumerate}

\begin{popfigure}
\begin{center}
\begin{tikzpicture}[node distance=0.55cm and 0.9cm,
bulle/.style={rounded corners=6pt, draw=popblue, fill=popblueL, thick, align=center, minimum height=1cm, text width=3.1cm, font=\small},
fleche/.style={-{Stealth[length=3mm]}, thick, popdark}]
\node[bulle, align=center] (g){\textbf{1. Greeting}\\ \emph{Good morning!}};
\node[bulle, right=of g, fill=poppurpleL, draw=poppurple, align=center] (p){\textbf{2. Purpose}\\ \emph{I'd like to apply for...}};
\node[bulle, right=of p, fill=poporangeL, draw=poporange, align=center] (q){\textbf{3. Questions}\\ \emph{Do you have...? How much...?}};
\node[bulle, below=of q, fill=popgreenL, draw=popgreen, align=center] (a){\textbf{4. Answers}\\ \emph{It's 75,000 francs.}};
\node[bulle, left=of a, fill=popgoldL, draw=popgold, align=center] (t){\textbf{5. Thanks \& Goodbye}\\ \emph{Thank you! You're welcome.}};
\draw[fleche] (g) -- (p);
\draw[fleche] (p) -- (q);
\draw[fleche] (q) -- (a);
\draw[fleche] (a) -- (t);
\end{tikzpicture}
\end{center}
\legende{Les cinq étapes d'un dialogue au guichet : de la salutation à la conclusion.}
\end{popfigure}

\begin{methode}[Comprendre un dialogue à l'écoute]
\begin{enumerate}
\item \textbf{Identifier QUI parle et OÙ} : écoutez les premiers mots (\eng{Good morning, can I help you?}) — ils révèlent le lieu (un guichet, un bureau) et les rôles (agent / usager).
\item \textbf{Repérer l'objet de la visite} : cherchez la phrase avec \eng{I'd like to...} ou \eng{I want to...} — elle donne le sujet du dialogue.
\item \textbf{Noter les chiffres} : prix (\eng{fee}), délais (\eng{two weeks}), nombre de documents (\eng{two photographs}). Ce sont presque toujours les réponses aux questions d'examen !
\item \textbf{Repérer la formule de conclusion} : \eng{Thank you}, \eng{You're welcome}, \eng{Goodbye} — elles signalent la fin de l'échange.
\end{enumerate}
\end{methode}

\courssub{Identification des formules fonctionnelles}

Relevons maintenant, dans le dialogue, les \cle{formules fonctionnelles} : ce sont des expressions toutes faites, à mémoriser comme des blocs, qui permettent de demander, donner une information, relancer et conclure.

\begin{center}
\begin{tabularx}{\linewidth}{|l|X|X|}\hline
\rowcolor{popblueL} \textbf{Fonction} & \textbf{English} & \textbf{Français} \\ \hline
Saluer / offrir son aide & Good morning, can I help you? & Bonjour, puis-je vous aider ? \\ \hline
Annoncer l'objet de la visite & I'd like to apply for a passport, please. & Je voudrais faire une demande de passeport, s'il vous plaît. \\ \hline
Demander une information & Do you have your birth certificate? & Avez-vous votre acte de naissance ? \\ \hline
Relancer / demander un complément & What else do I need? & Que me faut-il d'autre ? \\ \hline
Demander un prix & How much is the fee? & Combien coûtent les frais ? \\ \hline
Demander un délai & How long does it take? & Combien de temps cela prend-il ? \\ \hline
Donner une information chiffrée & It's 75,000 francs. / About two weeks. & C'est 75 000 francs. / Environ deux semaines. \\ \hline
Remercier & Thank you very much. & Merci beaucoup. \\ \hline
Répondre à un remerciement & You're welcome. & Je vous en prie. \\ \hline
\end{tabularx}
\end{center}

\begin{exemplebox}[Trois formules à retenir absolument]
\eng{What else do I need?} = « Que me faut-il d'autre ? » — \eng{else} (« d'autre ») se place \emph{après} le mot interrogatif : \eng{what else}, \eng{who else}, \eng{anything else}.

\eng{Here they are.} = « Les voici. » — on présente les documents à l'agent. Au singulier : \eng{Here it is.} (« Le voici »).

\eng{It takes about two weeks.} = « Cela prend environ deux semaines. » — le verbe \eng{to take} exprime la durée nécessaire : \eng{How long does it take? — It takes ten days.}
\end{exemplebox}

\begin{propriete}[Demander un prix et un délai en anglais]
\textbf{Pour un prix}, deux constructions sont correctes :
\begin{itemize}
\item \eng{How much is the fee?} (structure : \eng{How much} + verbe \eng{to be} + sujet) ;
\item \eng{How much does the passport cost?} (structure : \eng{How much} + auxiliaire \eng{does} + sujet + verbe \eng{cost}).
\end{itemize}
\textbf{Pour un délai} : \eng{How long does it take?} — littéralement « combien de temps cela prend-il ? ». La réponse utilise \eng{it takes} + durée : \eng{It takes about two weeks.}

Comparaison avec le français : là où le français dit « Combien coûte le passeport ? » avec inversion du verbe et du sujet, l'anglais exige l'auxiliaire \eng{do/does} : \eng{How much does the passport cost?}
\end{propriete}

\begin{attention}[L'erreur classique du francophone]
Ne dites jamais \emph{« How much costs the passport? »} : en anglais, on ne fait \textbf{pas d'inversion} directe du verbe et du sujet après \eng{How much} dans ce type de question. Il faut l'auxiliaire \eng{does} :

\eng{How much does the passport cost?} ou plus simplement \eng{How much is the fee?}

De même, évitez \emph{« How many time does it take? »} : « temps » (durée) est indénombrable en anglais, on dit \eng{How \textbf{long} does it take?} et non « how many time ».
\end{attention}

\courssub{Questions de compréhension orale}

Écoutez à nouveau le dialogue (ou relisez-le), puis répondez en phrases complètes :

\begin{enumerate}
\item Where does the dialogue take place?
\item What does the applicant want to do?
\item What documents does the officer mention?
\item How much is the fee?
\item How long does it take to get the passport?
\end{enumerate}

\begin{exoresolu}[Compréhension du dialogue]
\textbf{Énoncé.} Répondez en anglais, par une phrase complète : \emph{What documents does the officer mention? How much is the fee? How long does it take?}
\tcblower
\textbf{Solution détaillée.}
\begin{itemize}
\item Documents : \eng{The officer mentions the birth certificate, the ID card, two passport-sized photographs and the receipt of payment.} (« L'agent mentionne l'acte de naissance, la carte d'identité, deux photos d'identité et le reçu de paiement. ») — méthode : on reprend la question et on la transforme en affirmation ; on énumère avec \eng{and} avant le dernier élément.
\item Frais : \eng{The fee is 75,000 francs.} — on reprend la structure \eng{It's...} en nommant le sujet : \eng{The fee is...}
\item Délai : \eng{It takes about two weeks.} — on conserve la structure \eng{it takes} + durée.
\end{itemize}
\end{exoresolu}

\courssub{Lecture à voix haute et activité de substitution}

\begin{experience}[Lecture en binôme, puis substitution]
\textbf{Étape 1 — Lecture à voix haute.} En binôme, lisez le dialogue : un élève est l'\eng{Officer}, l'autre l'\eng{Applicant}. Soignez l'accentuation : \eng{Good \textbf{MOR}ning}, \eng{passport-sized \textbf{PHO}tographs}, et la montée d'intonation sur les questions fermées (\eng{yes/no}) : \eng{Do you have your birth certificate?} (voix qui monte à la fin).

\textbf{Étape 2 — Substitution.} Rejouez le même dialogue en remplaçant les informations entre crochets par vos propres données :
\begin{itemize}
\item documents : \eng{voter card}, \eng{school certificate}, \eng{four photographs} ;
\item montant : \eng{50,000 francs}, \eng{100,000 francs} ;
\item délai : \eng{one week}, \eng{ten days}, \eng{a month}.
\end{itemize}
Exemple de réplique transformée : \eng{How much is the fee? — It's 50,000 francs.} / \eng{How long does it take? — About ten days.}

\textbf{Étape 3 — Changement de rôles.} Inversez les rôles et recommencez avec de nouvelles données.
\end{experience}

\begin{exercice}[À vous de jouer]
Complétez ce mini-dialogue avec les formules du dialogue modèle :
\begin{enumerate}
\item Officer : Good morning, ..................... ?
\item Applicant : I'd like to ..................... a passport, please.
\item Officer : Do you have your ..................... and your ID card?
\item Applicant : Yes, here ..................... . What ..................... do I need?
\item Applicant : How ..................... is the fee? And how ..................... does it take?
\item Applicant : Thank you very much. — Officer : ..................... .
\end{enumerate}
\end{exercice}

\begin{corrige}[Exercice « À vous de jouer »]
1. \eng{can I help you} — 2. \eng{apply for} — 3. \eng{birth certificate} — 4. \eng{they are} ; \eng{else} — 5. \eng{much} ; \eng{long} — 6. \eng{You're welcome}.
\end{corrige}

Vous savez désormais écouter un dialogue administratif, en repérer les cinq étapes et en extraire les informations essentielles (documents, frais, délais). Dans la section suivante, nous irons plus loin avec un deuxième dialogue modèle : cette fois, le demandeur ne demande plus seulement des informations, il \emph{dépose réellement sa demande} de passeport au guichet.

\coursec{Dialogue modèle 2 : déposer sa demande}

Dans la section précédente, nous avons appris à \cle{demander des informations} au guichet : quels documents apporter, combien coûte le passeport, combien de temps il faut attendre. Nous voici maintenant une semaine plus tard : le demandeur revient au commissariat ou au centre d'état civil avec tous ses documents. Cette fois, il ne pose plus de questions : il \cle{dépose sa demande} et répond aux questions de l'agent. C'est une situation très fréquente au Cameroun, et elle exige de savoir donner ses informations personnelles, dire sa date de naissance et épeler son nom en anglais.

\courssub{Le dialogue modèle : à l'écoute et à l'observation}

Lisez d'abord le dialogue en silence, puis à voix haute, en jouant les deux rôles. Soulignez toutes les questions de l'agent et toutes les informations données par le demandeur.

\begin{textebox}[At the passport office — submitting the application]
\textbf{Officer:} Good afternoon. What's your full name, please?\\
\textbf{Applicant:} My name is Amina Mbarga.\\
\textbf{Officer:} Could you spell your surname?\\
\textbf{Applicant:} Yes, M-B-A-R-G-A.\\
\textbf{Officer:} When and where were you born?\\
\textbf{Applicant:} I was born on the third of April, two thousand and five, in Douala.\\
\textbf{Officer:} What's your nationality?\\
\textbf{Applicant:} I'm Cameroonian.\\
\textbf{Officer:} And what's your occupation?\\
\textbf{Applicant:} I'm a student at Government High School Biyem-Assi.\\
\textbf{Officer:} Are you married or single?\\
\textbf{Applicant:} I'm single.\\
\textbf{Officer:} Fine. Please sign here. Your passport will be ready on the twentieth.\\
\textbf{Applicant:} Thank you very much. Goodbye!\\
\textbf{Officer:} Goodbye, and don't forget your receipt.
\end{textebox}

\textbf{Glossaire}

\begin{center}
\begin{tabularx}{\linewidth}{|l|l|X|}
\hline
\rowcolor{popblueL} \textbf{English} & \textbf{Nature} & \textbf{Français} \\
\hline
full name & nom & nom complet (prénom + nom de famille) \\
\hline
surname & nom & nom de famille \\
\hline
to spell & verbe & épeler \\
\hline
to be born & verbe & naître, être né \\
\hline
nationality & nom & nationalité \\
\hline
occupation & nom & profession, métier \\
\hline
single & adjectif & célibataire \\
\hline
married & adjectif & marié(e) \\
\hline
to sign & verbe & signer \\
\hline
receipt & nom & récépissé, reçu (se prononce « ri-site », le \emph{p} est muet) \\
\hline
\end{tabularx}
\end{center}

\textbf{Questions de compréhension}

\begin{enumerate}
\item What is the applicant's full name?
\item When and where was she born?
\item What is her occupation?
\item When will her passport be ready?
\end{enumerate}

\begin{corrige}[Questions de compréhension]
1. \eng{Her full name is Amina Mbarga.}\\
2. \eng{She was born on the third of April, two thousand and five, in Douala.}\\
3. \eng{She is a student at Government High School Biyem-Assi.}\\
4. \eng{Her passport will be ready on the twentieth.}
\end{corrige}

\courssub{Donner ses informations personnelles}

Observons les questions de l'agent et les réponses d'Amina. Chaque question correspond à une rubrique du formulaire officiel.

\begin{center}
\begin{tabularx}{\linewidth}{|X|X|X|}
\hline
\rowcolor{popblueL} \textbf{Question de l'agent} & \textbf{Réponse possible} & \textbf{Français} \\
\hline
What's your full name? & My name is Amina Mbarga. & Quel est votre nom complet ? \\
\hline
What's your surname? & My surname is Mbarga. & Quel est votre nom de famille ? \\
\hline
When were you born? & I was born on the third of April, two thousand and five. & Quand êtes-vous né(e) ? \\
\hline
Where were you born? & I was born in Douala. & Où êtes-vous né(e) ? \\
\hline
What's your nationality? & I'm Cameroonian. & Quelle est votre nationalité ? \\
\hline
What's your address? & I live at 12, Avenue Kennedy, Yaoundé. & Quelle est votre adresse ? \\
\hline
What's your occupation? & I'm a student. / I'm a teacher. & Quelle est votre profession ? \\
\hline
Are you married or single? & I'm single. / I'm married. & Êtes-vous marié(e) ou célibataire ? \\
\hline
\end{tabularx}
\end{center}

\begin{definition}[Le vocabulaire de l'identité]
\eng{First name} (ou \eng{given name}) = prénom ; \eng{surname} (ou \eng{family name}, \eng{last name}) = nom de famille ; \eng{date of birth} = date de naissance ; \eng{place of birth} = lieu de naissance ; \eng{nationality} = nationalité ; \eng{address} = adresse ; \eng{occupation} = profession ; \eng{marital status} = situation matrimoniale (\eng{single} = célibataire, \eng{married} = marié, \eng{divorced} = divorcé, \eng{widowed} = veuf ou veuve).
\end{definition}

\begin{popfigure}
\begin{tikzpicture}[
box/.style={draw=popblue, fill=popblueL, rounded corners=2pt, align=center, font=\small, text width=3.4cm, minimum height=1.1cm},
ans/.style={draw=poporange, fill=poporangeL, rounded corners=2pt, align=center, font=\small, text width=4.6cm, minimum height=1.1cm},
lab/.style={font=\small\bfseries, text=popdark}]
\node[lab] at (1.7,6.6) {Personal information};
\node[box] (n1) at (1.7,5.6) {Name};
\node[ans] (a1) at (6.3,5.6) {Amina Mbarga};
\node[box] (n2) at (1.7,4.3) {Date of birth};
\node[ans] (a2) at (6.3,4.3) {3rd April 2005};
\node[box] (n3) at (1.7,3.0) {Place of birth};
\node[ans] (a3) at (6.3,3.0) {Douala, Cameroon};
\node[box] (n4) at (1.7,1.7) {Nationality};
\node[ans] (a4) at (6.3,1.7) {Cameroonian};
\node[box] (n5) at (1.7,0.4) {Address};
\node[ans] (a5) at (6.3,0.4) {12, Avenue Kennedy, Yaoundé};
\node[box] (n6) at (1.7,-0.9) {Occupation};
\node[ans] (a6) at (6.3,-0.9) {Student};
\draw[-{Stealth}, popline] (n1) -- (a1);
\draw[-{Stealth}, popline] (n2) -- (a2);
\draw[-{Stealth}, popline] (n3) -- (a3);
\draw[-{Stealth}, popline] (n4) -- (a4);
\draw[-{Stealth}, popline] (n5) -- (a5);
\draw[-{Stealth}, popline] (n6) -- (a6);
\end{tikzpicture}
\legende{Fiche « Personal information » remplie pour une élève camerounaise : chaque rubrique du formulaire correspond à une question de l'agent.}
\end{popfigure}

\courssub{« I was born » : parler de sa naissance}

\begin{attention}[Jamais « I born » ni « I am born »]
En français, on dit « je \textbf{suis né} » avec l'auxiliaire \emph{être}. En anglais, la naissance s'exprime avec le \textbf{passif au passé} : \eng{to be born} devient \eng{I was born} (et non \eng{I born}, \eng{I am born} ou \eng{I borned}).\\
\eng{I was born in Yaoundé.} = Je suis né à Yaoundé.\\
\eng{She was born in 2005.} = Elle est née en 2005.\\
\eng{They were born in Buea.} = Ils sont nés à Buea.\\
La question correspondante est : \eng{When were you born?} / \eng{Where were you born?}
\end{attention}

\begin{propriete}[Forme et emploi de « to be born »]
\begin{itemize}
\item \textbf{Forme} : sujet + \eng{was/were} + \eng{born} (+ complément de temps ou de lieu).
\item \eng{was} pour \eng{I, he, she, it} ; \eng{were} pour \eng{you, we, they}.
\item \textbf{Emploi} : uniquement pour la naissance d'une personne (ou d'un animal). On n'utilise jamais le présent : même un bébé né aujourd'hui \eng{was born today}.
\item \textbf{Prépositions} : \eng{born in} + ville, pays, mois ou année (\eng{born in Douala}, \eng{born in 2005}) ; \eng{born on} + date précise (\eng{born on the third of April}).
\end{itemize}
\end{propriete}

\begin{exemplebox}[Exemples traduits]
\eng{Could you spell your surname, please?} = Pourriez-vous épeler votre nom de famille, s'il vous plaît ?\\
\eng{I was born on the third of April.} = Je suis né le trois avril.\\
\eng{My brother was born in Bamenda in 2002.} = Mon frère est né à Bamenda en 2002.\\
\eng{Please sign here.} = Veuillez signer ici.\\
\eng{Your passport will be ready on the twentieth.} = Votre passeport sera prêt le vingt.
\end{exemplebox}

\courssub{Dire sa date de naissance en anglais britannique}

\begin{propriete}[La date à l'oral : « the fifteenth of May, nineteen ninety-eight »]
En anglais \textbf{britannique}, la date 15/05/1998 se dit : \eng{the fifteenth of May, nineteen ninety-eight} (« the fif-teenth of May, nine-teen nine-ty-eight »).\\
\textbf{Structure} : \eng{the} + nombre ordinal (\eng{first, second, third, fourth, fifth... twentieth, twenty-first, thirty-first}) + \eng{of} + mois + année.\\
L'\textbf{année} se lit en deux groupes de deux chiffres : 1998 = \eng{nineteen ninety-eight} ; 2005 = \eng{two thousand and five} ; 2010 = \eng{twenty ten}.\\
À l'écrit, on note : \eng{15th May 1998} ou \eng{15/05/1998} (jour/mois/année en Grande-Bretagne).
\end{propriete}

\begin{center}
\begin{tabular}{|l|l|l|}
\hline
\rowcolor{popblueL} \textbf{Date écrite} & \textbf{Date dite (britannique)} & \textbf{Français} \\
\hline
03/04/2005 & the third of April, two thousand and five & le 3 avril 2005 \\
\hline
21/07/2004 & the twenty-first of July, two thousand and four & le 21 juillet 2004 \\
\hline
01/01/2006 & the first of January, two thousand and six & le 1er janvier 2006 \\
\hline
30/11/1999 & the thirtieth of November, nineteen ninety-nine & le 30 novembre 1999 \\
\hline
\end{tabular}
\end{center}

\begin{savaistu}[« The third of April » ou « April third » ?]
En anglais \textbf{britannique}, la date se dit \eng{the third of April} et s'écrit 03/04 (jour puis mois). En anglais \textbf{américain}, on dit souvent \eng{April third} et on écrit 04/03 (mois puis jour) ! Résultat : la date 04/03 signifie le 4 mars à Londres mais le 3 avril à Washington. Sur les formulaires internationaux, méfiez-vous toujours de ce malentendu et, si possible, écrivez le mois en toutes lettres : \eng{3 April 2005}.
\end{savaistu}

\courssub{Épeler son nom : l'alphabet anglais en action}

Quand l'agent demande \eng{Could you spell your surname?}, il faut prononcer chaque lettre clairement. C'est une compétence indispensable au téléphone ou au guichet, car les noms camerounais (Ngo Bell, Etoundi, Tchoupo) sont inconnus des anglophones.

\begin{methode}[Comment épeler un nom en anglais]
\begin{enumerate}
\item Annoncez d'abord le nom entier : \eng{My surname is Ngo Bell.}
\item Ajoutez \eng{that's...} puis épelez lettre par lettre : \eng{that's N-G-O, B-E-L-L}.
\item Prononcez chaque lettre clairement, en marquant une petite pause entre les groupes pour les noms composés ou longs : \eng{E-T-O-U-N-D-I} peut se dire « E-T-O-U... N-D-I ».
\item Révisez les lettres pièges pour les francophones : A se dit « éï », E « i », I « aï », G « dji », J « djéï », H « éïtch », R « ar », W « da-beul-you », Y « ouaï ».
\end{enumerate}
\end{methode}

\begin{exemplebox}[L'alphabet anglais : les lettres à ne pas confondre]
A « éï » — E « i » — I « aï » — G « dji » — J « djéï » — H « éïtch » — K « kéï » — Q « kiou » — R « ar » — U « you » — V « vi » — W « da-beul-you » — X « èks » — Y « ouaï » — Z « zèd » (britannique) ou « zi » (américain).\\
\eng{My surname is Mbarga, that's M-B-A-R-G-A.} = Mon nom est Mbarga, ça s'écrit M-B-A-R-G-A.
\end{exemplebox}

\courssub{Signer, vérifier et conclure}

La fin de l'entretien est rituelle : l'agent demande une signature, annonce la date de retrait et remet un récépissé. Trois réflexes :

\begin{itemize}
\item \textbf{Signer} : \eng{Please sign here.} — répondez simplement \eng{Of course.} ou \eng{There you are.} en signant.
\item \textbf{Vérifier le récépissé} : contrôlez l'orthographe de votre nom et la date annoncée ; n'hésitez pas à demander \eng{Could you check the spelling of my name, please?}
\item \textbf{Noter la date de retrait} : \eng{Your passport will be ready on the twentieth.} — reformulez pour confirmer : \eng{So I come back on the twentieth. Thank you very much. Goodbye!}
\end{itemize}

\begin{exoresolu}[À vous de répondre à l'agent]
Complétez les réponses du demandeur avec vos propres informations (ou celles du personnage indiqué).\\
1. Officer: What's your full name, please? — Applicant: \ldots\\
2. Officer: Could you spell your surname? — Applicant: \ldots\\
3. Officer: When and where were you born? (personnage : né le 12 février 2005 à Bafoussam) — Applicant: \ldots\\
4. Officer: What's your occupation? (personnage : lycéen) — Applicant: \ldots\\
5. Officer: Please sign here. Your passport will be ready on the twenty-fifth. — Applicant: \ldots
\tcblower
\textbf{Solution détaillée}\\
1. \eng{My name is Jean-Paul Etoundi.} (ou votre propre nom : \eng{My name is + prénom + nom}.)\\
2. \eng{Yes, that's E-T-O-U-N-D-I.} (on épelle lettre par lettre, clairement, avec une pause éventuelle : « E-T-O-U... N-D-I ».)\\
3. \eng{I was born on the twelfth of February, two thousand and five, in Bafoussam.} (attention : \eng{I was born}, jamais \eng{I born} ; \eng{on} devant la date complète, \eng{in} devant la ville.)\\
4. \eng{I'm a student.} (ou \eng{I'm a high school student.})\\
5. \eng{Of course... So I come back on the twenty-fifth. Thank you very much. Goodbye!} (on signe, on reformule la date pour confirmer, on remercie et on prend congé.)
\end{exoresolu}

\begin{attention}[Les pièges des francophones au guichet]
\begin{itemize}
\item \eng{I born} ou \eng{I am born} : faux ! Dites toujours \eng{I was born}.
\item \eng{I have 17 years} : faux ! L'âge se dit avec \eng{to be} : \eng{I am seventeen (years old)}.
\item \eng{My name is called...} : faux ! Dites \eng{My name is...} ou \eng{I am called...}, mais ne mélangez pas les deux.
\item La date 05/06 : à l'oral britannique, dites \eng{the fifth of June} ; ne traduisez jamais mot à mot « le cinq juin » par \eng{the five June}.
\item \eng{sign} se prononce « saïne » (le \emph{g} est muet) ; \eng{signature} se prononce « si-gna-tcheur ».
\end{itemize}
\end{attention}

\begin{exercice}[Entraînement : votre fiche d'identité orale]
Préparez en une minute votre propre fiche, puis présentez-vous à votre voisin comme à un agent : nom complet, nom épelé, date et lieu de naissance, nationalité, adresse, profession. Votre voisin joue l'agent et pose les questions du tableau de la sous-section « Donner ses informations personnelles ». Inversez ensuite les rôles.
\end{exercice}

Vous savez maintenant déposer une demande de passeport de bout en bout : donner votre identité, dire votre date de naissance à la britannique, épeler votre nom et confirmer la date de retrait. Dans la section suivante, nous rassemblerons toutes ces formules dans une véritable \emph{boîte à outils fonctionnelle}, pour pouvoir réagir à toutes les situations du guichet.

\coursec{Boîte à outils : les formules fonctionnelles}

Dans les deux dialogues modèles que nous venons d'étudier — demander des informations puis déposer sa demande de passeport —, nous avons rencontré des phrases toutes faites qui reviennent sans cesse : \eng{Could you tell me...?}, \eng{You need...}, \eng{Sorry, I didn't catch that.} Ces phrases ne sont pas du vocabulaire isolé : ce sont des \cle{formules fonctionnelles}, c'est-à-dire des blocs de langue prêts à l'emploi qui remplissent une fonction précise dans la conversation (demander, donner, vérifier, conclure). Les mémoriser comme des unités entières, c'est se donner les moyens de parler vite, juste et poliment, sans avoir à construire chaque phrase mot à mot. Cette section les classe, les explique et vous apprend à les manier.

\courssub{Saluer et ouvrir la conversation}

Toute interaction à un guichet commence par une salutation et une annonce claire de l'objet de la visite. Observez les deux premiers échanges du dialogue modèle :

\begin{exemplebox}[Ouvrir la conversation]
\eng{Good morning, madam.} « Bonjour, madame. »\\
\eng{Good morning. Can I help you?} « Bonjour. Puis-je vous aider ? » (formule typique de l'employé au guichet)\\
\eng{I'd like to apply for a passport, please.} « Je voudrais faire une demande de passeport, s'il vous plaît. »\\
\eng{I need some information about passport renewal.} « J'ai besoin de renseignements sur le renouvellement du passeport. »
\end{exemplebox}

\begin{propriete}[I'd like... : la formule d'ouverture polie]
\eng{I'd like} est la contraction de \eng{I would like}. C'est l'équivalent poli du français « je voudrais ». Forme : \textbf{I'd like + nom} (\eng{I'd like an application form.}) ou \textbf{I'd like + to + base verbale} (\eng{I'd like to apply for a passport.}). On évite \eng{I want}, trop direct, qui sonne comme une exigence (« je veux »). Notez aussi le verbe \eng{to apply for} : « faire une demande de » ; la préposition \eng{for} est obligatoire (\eng{to apply for a passport}, \eng{to apply for a visa}).
\end{propriete}

\begin{attention}[Faux amis : « demander » et « assister »]
En français, « demander un passeport » se dit \eng{to apply for a passport} ou \eng{to ask for a passport}. Le verbe \eng{to demand} est un faux ami : il signifie « exiger » et serait très impoli au guichet. De même, \eng{to assist} ne veut pas dire « assister à » mais « aider » : \eng{Can I assist you?} signifie « Puis-je vous aider ? ».
\end{attention}

\courssub{Demander une information}

Les questions directes se construisent avec un mot interrogatif suivi de l'inversion (auxiliaire + sujet). Voici les cinq questions indispensables dans notre situation :

\begin{center}
\begin{tabularx}{\linewidth}{|X|X|X|}
\hline
\rowcolor{popblueL} English & Français & Remarque \\
\hline
\eng{What documents do I need?} & Quels documents me faut-il ? & \eng{need} + nom \\
\hline
\eng{How much is the fee?} & Combien coûte la taxe ? & \eng{how much} = prix (indénombrable) \\
\hline
\eng{How long does it take?} & Combien de temps cela prend-il ? & \eng{it takes} = « cela prend » \\
\hline
\eng{Where can I pay?} & Où puis-je payer ? & inversion de \eng{can} \\
\hline
\eng{When can I collect it?} & Quand puis-je le retirer ? & \eng{collect} = « retirer, récupérer » \\
\hline
\end{tabularx}
\end{center}

Pour un registre plus formel — recommandé face à un fonctionnaire —, on enrobe la question directe dans une formule indirecte :

\begin{exemplebox}[Registre formel]
\eng{Could you tell me what documents I need?} « Pourriez-vous me dire quels documents il me faut ? »\\
\eng{Could you repeat that, please?} « Pourriez-vous répéter, s'il vous plaît ? »\\
\eng{Would you mind spelling that?} « Cela vous dérangerait-il de l'épeler ? » (attention : \eng{mind} + verbe en \emph{-ing})\\
\eng{Could you tell me how long it takes?} « Pourriez-vous me dire combien de temps cela prend ? »
\end{exemplebox}

\begin{propriete}[La question indirecte : pas d'inversion]
Après \eng{Could you tell me}, \eng{Do you know}, \eng{I'd like to know}, la question devient une \cle{subordonnée interrogative indirecte} : on revient à l'ordre sujet + verbe, comme dans une phrase affirmative. Forme : \textbf{Could you tell me + mot interrogatif + sujet + verbe ?}\\
\eng{Where is the office?} (question directe, inversion) $\rightarrow$ \eng{Could you tell me where the office is?} (question indirecte, ordre normal). Si la question directe n'a pas de mot interrogatif, on utilise \eng{if} ou \eng{whether} : \eng{Could you tell me if the office is open on Saturdays?} « Pourriez-vous me dire si le bureau est ouvert le samedi ? »
\end{propriete}

\begin{attention}[Le piège classique du francophone]
Le français dit « Pouvez-vous me dire où est le bureau ? », avec inversion. En anglais, \eng{Could you tell me where is the office?} est une \textbf{erreur}. L'anglais interdit l'inversion dans la question indirecte : \eng{Could you tell me where the office is?} Autre exemple : \eng{Do you know when the office closes?} (et non \eng{when does the office close}).
\end{attention}

\courssub{Donner une information}

L'employé du service répond avec des formules simples et prévisibles. Apprenez-les aussi : dans un jeu de rôle, vous serez tantôt le demandeur, tantôt l'agent.

\begin{exemplebox}[Donner une information]
\eng{You need your birth certificate, a national ID card and two passport photos.} « Il vous faut votre acte de naissance, une carte nationale d'identité et deux photos d'identité. »\\
\eng{It costs 75,000 francs.} « Cela coûte 75 000 francs. »\\
\eng{It takes about two weeks.} « Cela prend environ deux semaines. »\\
\eng{You can pay at the cashier's office on the ground floor.} « Vous pouvez payer à la caisse au rez-de-chaussée. »\\
\eng{You can collect it on Monday.} « Vous pouvez le retirer lundi. »\\
\eng{First, you fill in this form; then you go to room 12.} « D'abord, vous remplissez ce formulaire ; ensuite, vous allez au bureau 12. »
\end{exemplebox}

Notez les verbes à la troisième personne du singulier avec \eng{it} : \eng{it cost\textbf{s}}, \eng{it take\textbf{s}}. Le \emph{s} de la troisième personne est l'une des erreurs les plus fréquentes des francophones à l'oral. Notez aussi le phrasal verb \eng{to fill in a form} (« remplir un formulaire ») : la particule \eng{in} est inséparable du sens.

\courssub{Relancer et vérifier : survivre à l'incompréhension}

Dans une vraie conversation, on ne comprend pas tout du premier coup. Un bon locuteur n'est pas celui qui comprend tout, mais celui qui sait \cle{relancer} sans paniquer. Observez :

\begin{exemplebox}[Relancer et vérifier]
\eng{Sorry, I didn't catch that.} « Pardon, je n'ai pas saisi. »\\
\eng{Do you mean the birth certificate?} « Vous voulez dire l'acte de naissance ? »\\
\eng{So, if I understand correctly, I come back in two weeks?} « Donc, si je comprends bien, je reviens dans deux semaines ? »\\
\eng{Could you say that again, please?} « Pourriez-vous redire cela, s'il vous plaît ? »\\
\eng{What do you mean by “certified copy”?} « Qu'entendez-vous par “copie certifiée” ? »
\end{exemplebox}

\begin{methode}[Relancer sans stress, en trois étapes]
\begin{enumerate}
\item \textbf{S'excuser brièvement} : \eng{Sorry,} / \eng{Excuse me,} — un seul mot suffit, pas de longues justifications.
\item \textbf{Nommer le problème} : \eng{I didn't catch the date.} / \eng{I don't understand.} / \eng{I'm not sure about the price.}
\item \textbf{Poser une question précise} : \eng{Could you repeat the date, please?} / \eng{Do you mean the original or a copy?}
\end{enumerate}
Cette stratégie en trois temps transforme un silence embarrassé en échange efficace. Elle est évaluée positivement à l'oral.
\end{methode}

\begin{attention}[« Pardon ? » ne se dit pas à la française]
Pour demander de répéter, on dit \eng{Sorry?} (avec l'intonation montante) ou \eng{Pardon?}. Le mot français « pardon ? » prononcé à la française est incompréhensible pour un anglophone : dites « so-ri ? » avec la voix qui monte. Évitez aussi \eng{What?} seul, très impoli.
\end{attention}

\courssub{Remercier et conclure}

Une conversation réussie se termine proprement : remerciement, réponse au remerciement, salutation finale.

\begin{exemplebox}[Conclure]
\eng{Thank you very much for your help.} « Merci beaucoup pour votre aide. »\\
\eng{You're welcome.} « Je vous en prie. / De rien. »\\
\eng{Have a nice day. Goodbye!} « Bonne journée. Au revoir ! »\\
\eng{Thanks. I'll come back on Monday, then.} « Merci. Je reviendrai lundi, alors. »
\end{exemplebox}

\begin{propriete}[L'échelle de politesse : plus c'est long et indirect, plus c'est poli]
En anglais, la politesse se mesure à la longueur et à l'indirection de la formule. Pour obtenir le formulaire, du plus sec au plus poli :\\
\eng{Give me the form.} (impératif sec, à éviter) $<$ \eng{Can I have the form?} (correct, courant) $<$ \eng{Could I have the form, please?} (poli) $<$ \eng{Would you mind giving me the form?} (très poli, formel).\\
Le conditionnel (\eng{could}, \eng{would}) adoucit toujours la demande. Au guichet comme à l'examen, visez au minimum le niveau \eng{Could I...?}.
\end{propriete}

\courssub{Tableau récapitulatif des formules}

\begin{popfigure}
\begin{tikzpicture}[node distance=0.25cm]
\tikzset{col/.style={rectangle, rounded corners=2pt, draw=popline, fill=popcream, text width=3.6cm, align=left, inner sep=4pt, font=\small},
head/.style={rectangle, rounded corners=2pt, text=white, text width=3.6cm, align=center, inner sep=4pt, font=\small\bfseries}}
\node[head, fill=popblue] (h1) {Asking};
\node[col, below=of h1, align=center] (c1){What documents do I need?\\ How much is the fee?\\ Could you tell me where the office is?};
\node[head, fill=popgreen, right=0.35cm of h1] (h2) {Giving information};
\node[col, below=of h2, align=center] (c2){You need your ID card.\\ It takes about two weeks.\\ You can collect it on Monday.};
\node[head, fill=poporange, right=0.35cm of h2] (h3) {Checking};
\node[col, below=of h3, align=center] (c3){Sorry, I didn't catch that.\\ Do you mean the birth certificate?\\ So, if I understand correctly...?};
\node[head, fill=poppurple, right=0.35cm of h3] (h4) {Closing};
\node[col, below=of h4, align=center] (c4){Thank you for your help.\\ You're welcome.\\ Have a nice day. Goodbye!};
\end{tikzpicture}
\legende{Les quatre fonctions de la conversation au guichet, avec trois formules clés chacune.}
\end{popfigure}

\begin{center}
\begin{tabularx}{\linewidth}{|X|X|}
\hline
\rowcolor{popblueL} Fonction & Formules à mémoriser \\
\hline
Saluer, ouvrir & \eng{Good morning. Can I help you? / I'd like to apply for a passport, please.} \\
\hline
Demander & \eng{What documents do I need? / How much is the fee? / How long does it take?} \\
\hline
Demander poliment & \eng{Could you tell me...? / Could you repeat that, please? / Would you mind spelling that?} \\
\hline
Donner une information & \eng{You need... / It costs... / It takes about two weeks. / You can collect it on Monday.} \\
\hline
Relancer, vérifier & \eng{Sorry, I didn't catch that. / Do you mean...? / So, if I understand correctly...?} \\
\hline
Remercier, conclure & \eng{Thank you very much for your help. / You're welcome. / Have a nice day.} \\
\hline
\end{tabularx}
\end{center}

\begin{exoresolu}[Remettre le dialogue dans l'ordre]
Énoncé — Voici les répliques d'un échange au service des passeports, dans le désordre. Numérotez-les de 1 à 6 pour reconstituer le dialogue.\\
a) \eng{It takes about two weeks. You can collect it on the 15th.}\\
b) \eng{Good morning. I'd like to apply for a passport, please.}\\
c) \eng{Thank you very much for your help. Goodbye!}\\
d) \eng{Good morning. Can I help you?}\\
e) \eng{Sorry, I didn't catch that. How long does it take?}\\
f) \eng{Certainly. You need your birth certificate and two photos.}
\tcblower
Solution détaillée — On suit la logique des fonctions : salutation de l'agent (d, 1), ouverture du demandeur (b, 2), information donnée par l'agent (f, 3), relance du demandeur qui n'a pas saisi (e, 4), réponse de l'agent (a, 5), conclusion (c, 6). Ordre correct : \textbf{1-d, 2-b, 3-f, 4-e, 5-a, 6-c}. Remarquez que la relance (e) contient une question directe avec inversion (\eng{How long does it take?}), car elle n'est pas introduite par \eng{Could you tell me}.
\end{exoresolu}

\begin{savaistu}[À l'oral du Bac, relancer rapporte des points]
Beaucoup de candidats croient qu'un bon oral est un long monologue récité sans faute. C'est faux : l'épreuve orale valorise l'\cle{aisance interactionnelle}, c'est-à-dire la capacité à réagir, échanger et gérer les imprévus. Un candidat qui dit calmement \eng{Sorry, could you say that again?} ou \eng{Do you mean...?} montre qu'il communique réellement — et cela rapporte plus de points qu'un texte appris par cœur. Les examinateurs préfèrent un dialogue vivant, même imparfait, à une récitation figée.
\end{savaistu}

Vous disposez maintenant de toutes les formules nécessaires, classées par fonction. Il reste à les faire vivre : la section suivante travaillera la prononciation, l'accentuation et l'intonation, avant les jeux de rôle où vous mettrez cette boîte à outils en pratique.

\coursec{Prononciation, accentuation et intonation}

Dans la section précédente, nous avons constitué notre boîte à outils de formules fonctionnelles : demander une information (\eng{Could you tell me...?}), donner une information (\eng{You will need...}), relancer (\eng{Sorry, what was that again?}) et conclure (\eng{Thank you for your help.}). Mais une formule correcte prononcée avec une mélodie française peut devenir incompréhensible, ou pire, paraître impolie. Un agent de l'\eng{immigration office} de Yaoundé ou de Douala jugera votre anglais autant sur votre prononciation que sur votre grammaire. Cette section vous apprend donc à \cle{accentuer} correctement les mots clés, à prononcer les sons difficiles comme le \eng{th}, et à faire \cle{monter} ou \cle{descendre} la voix au bon moment.

\courssub{Observer d'abord : deux questions, deux mélodies}

\begin{textebox}[At the passport office — extrait]
\textbf{Officer:} Good morning. Can I help you? ↗\\
\textbf{Applicant:} Yes, please. I would like to apply for a PASSport. ↘\\
\textbf{Officer:} Do you have your birth cerTIFicate? ↗\\
\textbf{Applicant:} Yes, here it is. ↘ And what other documents do I need? ↘\\
\textbf{Officer:} You need two passport PHOtographs and your ID card. ↘\\
\textbf{Applicant:} How much is the fee? ↘\\
\textbf{Officer:} It is seventy-five thousand francs. ↘ Here is your reCEIPT. ↘\\
\textbf{Applicant:} Could you repeat that, please? ↗\\
\textbf{Officer:} Of course. Seventy-five thousand francs. ↘
\end{textebox}

Écoutez votre professeur lire ce dialogue et observez : la voix ne reste jamais plate. Elle \textbf{monte} à la fin de certaines questions et \textbf{descend} à la fin d'autres phrases. Certains mots sont prononcés plus fort que d'autres : \eng{PASSport}, \eng{cerTIFicate}, \eng{PHOtographs}, \eng{reCEIPT}. Ce n'est pas un hasard : c'est un système.

\courssub{L'accentuation des mots : frapper la bonne syllabe}

\begin{definition}[L'accent tonique (word stress)]
En anglais, chaque mot de plus d'une syllabe possède une syllabe \cle{accentuée}, prononcée plus fort, plus longtemps et plus haut que les autres. Contrairement au français, où l'accent tombe presque toujours sur la dernière syllabe du groupe (\emph{un passePORT, un certifiCAT}), l'anglais place l'accent selon des habitudes propres à chaque mot : il faut l'apprendre en même temps que le mot.
\end{definition}

Voici les mots clés de notre situation de communication, avec leur accentuation (la syllabe accentuée est en majuscules) et leur prononciation figurée en lettres françaises :

\begin{center}
\begin{tabularx}{\linewidth}{|l|X|X|X|}\hline
\rowcolor{popblueL} \textbf{Mot anglais} & \textbf{Accentuation} & \textbf{Prononciation figurée} & \textbf{Traduction} \\ \hline
passport & \cle{PASS}port & « passe-porte » & passeport \\ \hline
application & appli\cle{CA}tion & « a-pli-KÉ-cheune » & demande, candidature \\ \hline
certificate & cer\cle{TIF}icate & « seur-TI-fi-kète » & acte, certificat \\ \hline
photograph & \cle{PHO}tograph & « FÔ-teu-grafe » & photographie \\ \hline
receipt & re\cle{CEIPT} & « ri-SSITE » (p muet) & reçu, quittance \\ \hline
embassy & \cle{EM}bassy & « EM-be-si » & ambassade \\ \hline
officer & \cle{OF}ficer & « O-fi-seur » & agent, fonctionnaire \\ \hline
document & \cle{DO}cument & « DO-kieu-mente » & document \\ \hline
\end{tabularx}
\end{center}

\begin{attention}[Attention : certificate ne s'accentue pas comme certificat]
Le français dit \emph{un certifiCAT} (accent final) ; l'anglais dit \eng{cerTIFicate} : l'accent frappe la syllabe \textbf{tif}, ni la première ni la dernière. C'est un bon exemple de mot qu'il faut mémoriser avec sa prononciation : \eng{a birth cerTIFicate} « un acte de naissance ».
\end{attention}

\begin{propriete}[Règle d'or : les noms composés s'accentuent sur le premier élément]
Quand deux mots s'unissent pour former un nom composé, l'accent tonique principal tombe presque toujours sur le \cle{premier} élément : \eng{PASSport} (pass + port), \eng{BIRTH certificate}, \eng{ID card} (« aï-DII card »), \eng{APplication form}, \eng{PASSport photograph}. C'est l'inverse du français, qui accentue la fin : \emph{un certificat de naisSANCE}.
\end{propriete}

\begin{attention}[Le piège du « p » muet dans receipt]
Le mot \eng{receipt} s'écrit avec un \textbf{p} mais se prononce « rissite » : le \textbf{p} est \cle{muet} (comme dans \eng{psychology}, « saï-ko-lo-dji »). Ne dites jamais « ré-sèpte ». De même, \eng{receipt} est un faux ami partiel : il signifie « reçu, quittance », pas « recette » (qui se dit \eng{recipe}, « ré-si-pi »).
\end{attention}

\courssub{Les sons difficiles pour les francophones}

\begin{attention}[Le « th » anglais n'existe pas en français]
Le son \eng{th} est la difficulté numéro un des francophones. La technique : \cle{placez le bout de la langue entre les dents} et soufflez. Il existe deux variantes :
\begin{itemize}
\item le \textbf{th sourd} (sans vibration des cordes vocales) : \eng{birth} (« beurth »), \eng{third} (« theurd »), \eng{three}, \eng{month} ;
\item le \textbf{th sonore} (avec vibration, comme un « z » soufflé entre les dents) : \eng{the} (« ze »), \eng{that} (« zate »), \eng{this}, \eng{mother}.
\end{itemize}
Ne dites jamais « birz » pour \eng{birth} ni « te » pour \eng{the} : un Anglophone entendrait un autre mot ou ne comprendrait pas.
\end{attention}

\begin{exemplebox}[Mots difficiles de la leçon, prononciation figurée et traduction]
\begin{itemize}
\item \eng{receipt} = « rissite » (p muet) = le reçu. \eng{Keep your receipt.} « Gardez votre reçu. »
\item \eng{certificate} = « seur-TI-fi-kète » = l'acte, le certificat. \eng{I need a birth certificate.} « J'ai besoin d'un acte de naissance. »
\item \eng{officer} = « O-fi-seur » = l'agent. \eng{The officer checks my documents.} « L'agent vérifie mes documents. »
\item \eng{fee} = « fii » (i long) = les frais. \eng{How much is the fee?} « Combien coûtent les frais ? »
\item \eng{form} = « fome » (o bref, r non roulé) = le formulaire. \eng{Fill in this form, please.} « Remplissez ce formulaire, s'il vous plaît. »
\item \eng{third} = « theurd » = troisième. \eng{Go to the third window.} « Allez au troisième guichet. »
\end{itemize}
\end{exemplebox}

\courssub{L'intonation : la mélodie de la question}

\begin{propriete}[Les deux mélodies fondamentales de l'anglais]
\begin{enumerate}
\item \textbf{Intonation montante ↗} : pour les questions fermées (Yes/No questions), celles qui commencent par un auxiliaire (\eng{Do you...? Can I...? Have you...?}) et pour les demandes polies. La voix monte sur le dernier mot accentué.
\item \textbf{Intonation descendante ↘} : pour les questions ouvertes en \eng{Wh-} (\eng{What, Where, When, How much...}) et pour toutes les phrases affirmatives. La voix descend sur le dernier mot accentué.
\end{enumerate}
\end{propriete}

\begin{popfigure}
\begin{tikzpicture}[line width=1.6pt]
% Flèche montante
\draw[->, popgreen] (0,0.2) .. controls (1.2,0.2) and (2.2,0.5) .. (3.4,1.8);
\node[align=center, text=popgreen] at (1.7,-0.5) {\textbf{Yes/No questions ↗}};
\node[align=center, text=popdark, font=\small] at (1.7,-1.2) {Do you have your ID card? ↗};
% Flèche descendante
\draw[->, poporange] (6.2,1.8) .. controls (7.4,1.8) and (8.4,1.4) .. (9.6,0.2);
\node[align=center, text=poporange] at (7.9,-0.5) {\textbf{Wh- questions et affirmations ↘}};
\node[align=center, text=popdark, font=\small] at (7.9,-1.2) {What documents do I need? ↘};
\end{tikzpicture}
\legende{La voix monte quand on attend « yes » ou « no » ; elle descend quand on attend une information précise ou quand on affirme.}
\end{popfigure}

\begin{methode}[Choisir la bonne mélodie en trois questions]
\begin{enumerate}
\item \textbf{Ma phrase est-elle une question ?} Si non, la voix descend ↘ : \eng{It takes two weeks. ↘} « Cela prend deux semaines. »
\item \textbf{Si c'est une question, commence-t-elle par un mot interrogatif} (\eng{who, what, where, when, how much, how long}) ? Alors la voix descend ↘ : \eng{How much is the fee? ↘}
\item \textbf{Commence-t-elle par un auxiliaire} (\eng{do, can, could, have, is, will}) ? Alors la voix monte ↗, car j'attends \eng{yes} ou \eng{no} : \eng{Could you repeat that, please? ↗}
\end{enumerate}
Astuce de politesse : pour paraître courtois, utilisez une voix douce avec une montée légère : \eng{Could you help me, please? ↗} Une voix descendante et sèche sur une demande sonne comme un ordre.
\end{methode}

\begin{center}
\begin{tabularx}{\linewidth}{|X|X|X|}\hline
\rowcolor{popblueL} \textbf{Français (mélodie souvent plate)} & \textbf{English (mélodie marquée)} & \textbf{Type} \\ \hline
Vous avez votre carte d'identité ? & Do you have your ID card? ↗ & Yes/No → montante \\ \hline
De quels documents ai-je besoin ? & What documents do I need? ↘ & Wh- → descendante \\ \hline
Cela prend deux semaines. & It takes two weeks. ↘ & Affirmation → descendante \\ \hline
Pouvez-vous répéter, s'il vous plaît ? & Could you repeat that, please? ↗ & Demande polie → montante \\ \hline
Où est le bureau des passeports ? & Where is the passport office? ↘ & Wh- → descendante \\ \hline
\end{tabularx}
\end{center}

\begin{savaistu}[Pourquoi PASSport et non passPORT ?]
La plupart des noms composés anglais s'accentuent sur leur \cle{premier} élément : \eng{PASSport}, \eng{BIRTH certificate}, \eng{ID card}, \eng{APplication form}, \eng{FOOTball}, \eng{CLASSroom}. Cette habitude est si forte que les Anglophones reconnaissent un nom composé rien qu'à son accent. À l'inverse, les verbes à particule gardent souvent l'accent sur le second élément : \eng{to fill IN a form}. Le mot \eng{passport} vient d'ailleurs du français ancien \emph{passeport} (« passer le port »), emprunté au XVIe siècle — mais l'anglais a déplacé l'accent sur la première syllabe !
\end{savaistu}

\begin{exoresolu}[Trouve la bonne mélodie]
Énoncé. Pour chaque phrase, indiquez si l'intonation finale monte (↗) ou descend (↘), puis justifiez. a) \eng{Can I apply for a passport here?} b) \eng{When will my passport be ready?} c) \eng{You must pay the fee first.} d) \eng{Do you accept cash?} e) \eng{How long does it take?}
\tcblower
Solution détaillée. a) ↗ : la question commence par l'auxiliaire \eng{can}, on attend \eng{yes} ou \eng{no}. b) ↘ : question en \eng{Wh-} (\eng{when}), on attend une information précise (une date). c) ↘ : phrase affirmative, la voix descend sur \eng{first}. d) ↗ : question fermée commençant par \eng{do}. e) ↘ : question en \eng{Wh-} (\eng{how long}). Règle à retenir : auxiliaire en tête → la voix monte ; mot interrogatif en tête → la voix descend.
\end{exoresolu}

\courssub{Entraînement guidé : de la chorale au binôme}

\begin{experience}[Entraînement à la prononciation et à l'intonation]
\begin{enumerate}
\item \textbf{Répétition chorale.} Le professeur lit chaque phrase clé ; la classe répète en chœur en suivant la mélodie avec la main : la main monte ↗ pour \eng{Do you have your ID card? ↗}, la main descend ↘ pour \eng{It takes two weeks. ↘}.
\item \textbf{Travail sur les sons.} Répétez trois fois en chœur : \eng{birth, third, the, that, receipt, fee, form}. Vérifiez que la langue passe bien entre les dents pour \eng{th}.
\item \textbf{En binôme.} L'élève A lit une phrase du dialogue modèle ; l'élève B lève la main si la mélodie monte, la baisse si elle descend. Puis échangez les rôles.
\item \textbf{Défi final.} En binôme, rejouez le dialogue complet au guichet des passeports en soignant l'accentuation (\eng{PASSport, cerTIFicate, reCEIPT}) et les deux mélodies. La classe écoute et applaudit le binôme le plus mélodieux.
\end{enumerate}
\end{experience}

\begin{aretenir}[L'essentiel de la section]
\begin{itemize}
\item Chaque mot anglais a sa syllabe accentuée : \eng{PASSport, appliCAtion, cerTIFicate, PHOtograph, reCEIPT, EMbassy, OFficer}. Les noms composés s'accentuent sur le premier élément.
\item Le \eng{th} se prononce la langue entre les dents : sourd dans \eng{birth, third}, sonore dans \eng{the, that}. Le \eng{p} de \eng{receipt} est muet (« rissite »).
\item La voix \textbf{monte} ↗ pour les questions Yes/No et les demandes polies ; elle \textbf{descend} ↘ pour les questions en \eng{Wh-} et les affirmations.
\item Une bonne mélodie rend votre anglais plus clair et plus poli : entraînez-vous avec des gestes de la main.
\end{itemize}
\end{aretenir}

Votre prononciation est maintenant prête : vous savez frapper la bonne syllabe, souffler le \eng{th} et faire danser votre voix. Il ne reste plus qu'à mettre tout en scène dans la dernière section : les jeux de rôle guidés et la production libre au guichet des passeports.

\coursec{Jeux de rôle guidés et production libre}

Maintenant que vous maîtrisez la prononciation des mots clés (\eng{passport}, \eng{application}, \eng{expired}) et que vous savez placer l'accent et l'intonation montante des questions, il est temps de \cle{passer à l'action}. Parler, c'est comme jouer au football : on peut connaître toutes les règles, mais seul l'entraînement sur le terrain rend performant. Dans cette dernière étape, vous allez jouer de vraies scènes de guichet, d'abord très guidées, puis de plus en plus libres, jusqu'à improviser un dialogue complet comme vous pourriez le faire un jour à la délégation de l'immigration de Yaoundé ou dans une ambassade.

\courssub{Jeu de rôle 1 (guidé) : At the passport office}

Pour ce premier jeu de rôle, la classe travaille en \cle{binômes}. Un élève joue l'\eng{applicant} (le demandeur), l'autre joue l'\eng{officer} (l'agent du guichet). Chacun reçoit une fiche de rôle : c'est votre « scénario secret ». Ne montrez jamais votre fiche à votre partenaire !

\begin{textebox}[Role card A — Applicant]
You want a new passport. Go to the passport office and ask about:
\begin{itemize}
\item the documents you need;
\item the fee (how much does it cost?);
\item the processing time (how long does it take?);
\item the collection day (when and where can you collect it?).
\end{itemize}
You must ask at least four questions. Start with a greeting and end politely.
\end{textebox}

\begin{textebox}[Role card B — Officer]
You work at the passport office. Here is your information:
\begin{itemize}
\item Documents: birth certificate, ID card, 2 passport photos.
\item Fee: 75,000 FCFA.
\item Processing time: 2 weeks.
\item Collection: Monday to Friday, 10 a.m. to 3 p.m.
\end{itemize}
Answer the applicant's questions clearly. Greet the customer first and close the conversation politely.
\end{textebox}

\begin{exemplebox}[Un début de dialogue possible]
\eng{Officer: Good morning! Next, please! How can I help you?}\\
« Bonjour ! Au suivant, s'il vous plaît ! Comment puis-je vous aider ? »\\
\eng{Applicant: Good morning. I would like to apply for a passport. What documents do I need?}\\
« Bonjour. Je voudrais faire une demande de passeport. Quels documents me faut-il ? »\\
\eng{Officer: You need your birth certificate, your ID card and two passport photos.}\\
« Il vous faut votre acte de naissance, votre carte d'identité et deux photos d'identité. »
\end{exemplebox}

\begin{methode}[Réussir un jeu de rôle en quatre étapes]
\begin{enumerate}
\item \textbf{Préparer} : lisez votre fiche de rôle et écrivez au brouillon au moins trois questions (pour le demandeur) ou trois réponses (pour l'agent). Prévoyez aussi une formule de salutation et une formule de conclusion.
\item \textbf{Improviser} : pendant le dialogue, posez la fiche et parlez \emph{sans lire}. Regardez votre partenaire dans les yeux. Quelques hésitations sont normales : un vrai locuteur hésite aussi !
\item \textbf{Écouter pour réagir} : n'attendez pas simplement « votre tour de parole ». Écoutez vraiment la réponse pour enchaîner : \eng{Two weeks? That's quite long!} (« Deux semaines ? C'est plutôt long ! ») ou \eng{I see, thank you. And how much is the fee?} (« Je vois, merci. Et combien coûte la demande ? »).
\item \textbf{Conclure poliment} : terminez toujours par une formule de clôture : \eng{Thank you very much for your help. Goodbye!} / \eng{You're welcome. Have a nice day!} (« Merci beaucoup pour votre aide. Au revoir ! » / « Je vous en prie. Bonne journée ! »).
\end{enumerate}
\end{methode}

\courssub{Jeu de rôle 2 (semi-guidé) : At the Nigerian embassy}

Deuxième étape : on retire les phrases toutes faites. Vous ne recevez plus que les \cle{étapes} du dialogue ; à vous de trouver les mots. La situation : vous êtes à l'ambassade du Nigeria à Yaoundé et vous demandez un \eng{visa} pour rendre visite à votre oncle à Lagos.

\begin{experience}[Applying for a visa — dialogue skeleton]
Respectez l'ordre des étapes suivantes, mais inventez vos propres phrases :
\begin{enumerate}
\item The officer greets the applicant. (\eng{Good morning, how can I help you?})
\item The applicant states his or her purpose. (\eng{I would like to apply for a visa to Nigeria.})
\item The applicant asks about the required documents.
\item The officer gives the list of documents and the fee.
\item The applicant asks about the processing time and asks for clarification. (\eng{Sorry, did you say two weeks or two months?})
\item The officer repeats the information.
\item The applicant thanks the officer; both say goodbye.
\end{enumerate}
\end{experience}

\begin{exemplebox}[Demander une clarification — un réflexe indispensable]
\eng{Sorry, could you repeat that, please?} « Pardon, pourriez-vous répéter, s'il vous plaît ? »\\
\eng{Did you say fifteen or fifty?} « Vous avez dit quinze ou cinquante ? »\\
\eng{What do you mean by “supporting documents”?} « Qu'entendez-vous par “documents justificatifs” ? »\\
Ces formules de \cle{vérification de la compréhension} sont la marque d'un vrai communicant : mieux vaut faire répéter que repartir avec une fausse information.
\end{exemplebox}

\courssub{Jeu de rôle 3 (libre) : The expired passport}

Dernier niveau : la production libre. Plus de fiche détaillée, seulement une situation à problème. Votre passeport a expiré et vous devez le renouveler de toute urgence car vous partez en voyage scolaire dans dix jours. Mauvaise surprise : au guichet, vous réalisez que vous avez \cle{oublié un document}. Il va falloir \emph{négocier} une solution avec l'agent.

\begin{experience}[Renewing an expired passport — free role play]
\begin{itemize}
\item \textbf{Applicant} : your passport has expired. You forgot your birth certificate at home. Explain the problem, apologise, and negotiate: can you bring the document tomorrow? Can the process be accelerated?
\item \textbf{Officer} : the rules are strict, but you can be helpful. Accept or refuse, explain why, and propose a solution (come back tomorrow, pay an express fee, etc.).
\end{itemize}
\end{experience}

\begin{exemplebox}[Des phrases utiles pour négocier]
\eng{I'm afraid my passport has expired.} « Je crains que mon passeport n'ait expiré. »\\
\eng{I'm terribly sorry, I left my birth certificate at home.} « Je suis vraiment désolé(e), j'ai laissé mon acte de naissance à la maison. »\\
\eng{Is there any way to speed up the process?} « Y a-t-il un moyen d'accélérer la procédure ? »\\
\eng{Could I bring the missing document tomorrow morning?} « Pourrais-je apporter le document manquant demain matin ? »\\
\eng{That would be very kind of you.} « Ce serait très aimable de votre part. »
\end{exemplebox}

\begin{attention}[« Expired » : prononciation et grammaire]
Le mot \eng{expired} se prononce « èk-spaïeud » : le \eng{-ed} final se prononce bien « d » (et non « id »), et l'accent tombe sur la deuxième syllabe : ex-\textbf{PI}-red.\\
Côté grammaire, pour un résultat présent, on utilise le \cle{present perfect} : \eng{My passport has expired.} (« Mon passeport a expiré » — il est invalide maintenant). Ne dites pas \eng{My passport is expired} : le francophone qui traduit mot à mot « mon passeport est expiré » tombe dans le piège, car on exprime ici un événement dont le résultat est visible maintenant, d'où le present perfect. De même : \eng{My visa has run out.} « Mon visa a expiré. »
\end{attention}

\courssub{La checklist du dialogue réussi}

Avant de passer à l'auto-évaluation, visualisons les six étapes d'un dialogue de guichet réussi. Cette « rampe de lancement » doit devenir un réflexe.

\begin{popfigure}
\begin{center}
\begin{tikzpicture}[node distance=0.5cm,
  etape/.style={rectangle, rounded corners=3pt, draw=popblue, fill=popblueL, text=popdark, align=center, minimum width=4.6cm, minimum height=0.85cm, font=\small},
  fleche/.style={-{Stealth[length=2.5mm]}, thick, poporange}]
\node[etape, align=center] (n1){\textbf{1. Greet}\\ Good morning!};
\node[etape, below=of n1, align=center] (n2){\textbf{2. State purpose}\\ I would like to apply for\ldots};
\node[etape, below=of n2, align=center] (n3){\textbf{3. Ask 3+ questions}\\ What documents\ldots? How much\ldots?};
\node[etape, below=of n3, align=center] (n4){\textbf{4. Give personal information}\\ My name is\ldots I was born in\ldots};
\node[etape, below=of n4, align=center] (n5){\textbf{5. Check understanding}\\ Sorry, did you say\ldots?};
\node[etape, below=of n5, draw=popgreen, fill=popgreenL, align=center] (n6){\textbf{6. Thank \& close}\\ Thank you for your help. Goodbye!};
\draw[fleche] (n1) -- (n2);
\draw[fleche] (n2) -- (n3);
\draw[fleche] (n3) -- (n4);
\draw[fleche] (n4) -- (n5);
\draw[fleche] (n5) -- (n6);
\end{tikzpicture}
\end{center}
\legende{Checklist du dialogue réussi : les six étapes à suivre dans l'ordre, de la salutation à la clôture polie.}
\end{popfigure}

\courssub{Grille d'auto-évaluation}

Après chaque jeu de rôle, prenez une minute pour vous évaluer honnêtement. Cochez ce que vous avez réellement réussi à faire. Cette grille reprend les six étapes de la checklist : si une case reste vide, c'est votre objectif pour la prochaine prise de parole.

\begin{center}
\begin{tabularx}{\linewidth}{|X|l|l|l|}
\hline
\rowcolor{popblueL} \textbf{Critère} & \textbf{Oui} & \textbf{Un peu} & \textbf{Pas encore} \\
\hline
J'ai salué mon interlocuteur. (\eng{Good morning!}) & & & \\
\hline
J'ai annoncé clairement l'objet de ma visite. (\eng{I would like to\ldots}) & & & \\
\hline
J'ai posé des questions claires et bien formées. & & & \\
\hline
J'ai donné mes informations personnelles correctement. & & & \\
\hline
J'ai relancé ou fait répéter quand je n'ai pas compris. & & & \\
\hline
J'ai remercié et conclu poliment. (\eng{Thank you. Goodbye!}) & & & \\
\hline
\end{tabularx}
\end{center}

\begin{exoresolu}[Mettre en ordre un dialogue]
Voici un dialogue au guichet dont les répliques sont mélangées. Remettez-les dans l'ordre logique, puis jouez-le avec un camarade.\\
a. \eng{You need your ID card, two photos and the completed form.}\\
b. \eng{Good morning! How can I help you?}\\
c. \eng{Thank you very much. Goodbye!}\\
d. \eng{Good morning. I would like to apply for a passport. What documents do I need?}\\
e. \eng{It costs 75,000 FCFA and it takes two weeks.}\\
f. \eng{And how much is the fee, please? How long does it take?}
\tcblower
Ordre correct : \textbf{b → d → a → f → e → c}.\\
Justification : l'agent salue d'abord (b) ; le demandeur annonce son but et pose sa première question (d) ; l'agent donne la liste des documents (a) ; le demandeur enchaîne sur le tarif et le délai (f) ; l'agent répond (e) ; le demandeur remercie et conclut (c). On retrouve exactement la checklist : saluer → annoncer → questionner → informer → vérifier → conclure.
\end{exoresolu}

\courssub{Passage devant la classe et feedback collectif}

Le professeur invite ensuite les binômes les plus à l'aise à jouer leur scène \cle{devant toute la classe}. Le public n'est pas passif : chaque élève écoute avec la grille d'évaluation sous les yeux et note deux points forts et un point à améliorer.

\begin{experience}[Feedback collectif en deux temps]
\begin{enumerate}
\item \textbf{Prononciation} : la classe signale les mots bien prononcés et ceux à retravailler — accent sur la bonne syllabe (\eng{PASS-port}, ex-\textbf{PI}-red), intonation montante dans les questions fermées (\eng{Can I collect it on Friday?} ↗) et descendante dans les questions en \eng{wh-} (\eng{What documents do I need?} ↘).
\item \textbf{Politesse} : a-t-on entendu \eng{please}, \eng{thank you}, \eng{I'm afraid\ldots}, \eng{Could you\ldots?} ? En anglais, la politesse n'est pas un luxe : un \eng{Give me the form!} sec passerait pour de la grossièreté, alors qu'en français « donnez-moi le formulaire » peut rester neutre.
\end{enumerate}
\end{experience}

\begin{savaistu}[“Next, please!” : la formule du guichet]
Dans les pays anglophones, on fait souvent la queue devant le \eng{counter} (le guichet) — à la banque, à la poste, à l'aéroport. Quand c'est votre tour, l'agent lance : \eng{Next, please!} (« Au suivant, s'il vous plaît ! »), parfois suivi de \eng{Step forward, please.} (« Avancez, s'il vous plaît. »). Reconnaître ces formules à l'oreille est indispensable : si vous jouez le rôle de l'agent, utilisez-les pour ouvrir le dialogue de façon très authentique. À Londres comme à Lagos, la file d'attente (\eng{queue} en anglais britannique, \eng{line} en anglais américain) est une institution : on ne « grille » jamais son tour !
\end{savaistu}

\courssub{Complément Bac : raconter une démarche au passé}

\begin{aretenir}[Un atout pour l'oral du Baccalauréat]
À l'oral du Bac, on peut vous demander de raconter une expérience personnelle. Savoir décrire une \cle{démarche administrative au passé} est alors précieux. On utilise le \emph{prétérit simple} pour les actions accomplies, enchaînées par des marqueurs chronologiques : \eng{first}, \eng{then}, \eng{after that}, \eng{finally}.
\end{aretenir}

\begin{exemplebox}[Raconter sa démarche au passé]
\eng{First, I filled in the application form, then I paid the fee at the counter. After that, I handed in my documents, and finally, I collected my new passport two weeks later.}\\
« D'abord, j'ai rempli le formulaire de demande, puis j'ai payé les frais au guichet. Ensuite, j'ai déposé mes documents, et enfin, j'ai retiré mon nouveau passeport deux semaines plus tard. »
\end{exemplebox}

\begin{center}
\begin{tabularx}{\linewidth}{|X|X|}
\hline
\rowcolor{popblueL} \textbf{Marqueur chronologique} & \textbf{Exemple au prétérit} \\
\hline
\eng{First,} (d'abord) & \eng{First, I went to the passport office.} \\
\hline
\eng{Then,} (puis) & \eng{Then, I filled in the form.} \\
\hline
\eng{After that,} (ensuite) & \eng{After that, I paid 75,000 FCFA.} \\
\hline
\eng{Finally,} (enfin) & \eng{Finally, I collected my passport.} \\
\hline
\end{tabularx}
\end{center}

\begin{attention}[Attention aux verbes irréguliers au prétérit]
Le récit au passé exige les formes correctes des verbes irréguliers : \eng{go → went} (j'allai), \eng{pay → paid} (je payai — orthographe : \emph{paid}, pas « payed »), \eng{take → took} (cela prit), \eng{bring → brought} (j'apportai), \eng{forget → forgot} (j'oubliai). Une erreur fréquente : \eng{I bringed my documents} — non, c'est \eng{I brought my documents}. En revanche, les verbes réguliers prennent simplement \eng{-ed} : \eng{fill → filled}, \eng{collect → collected}, \eng{apply → applied} (le \eng{y} précédé d'une consonne devient \eng{i}).
\end{attention}

\begin{exercice}[À vous de jouer]
\begin{enumerate}
\item En binôme, rejouez le jeu de rôle 3 en inversant les rôles, puis remplissez chacun votre grille d'auto-évaluation.
\item Rédigez au brouillon cinq phrases au prétérit racontant votre dernière démarche administrative (réelle ou imaginaire), en utilisant au moins trois marqueurs chronologiques et deux verbes irréguliers.
\item Préparez-vous à raconter cette expérience oralement en trente secondes, sans lire.
\end{enumerate}
\end{exercice}

\begin{corrige}[Exercice — indications]
1. Veillez à ce que chaque partenaire ait joué les deux rôles ; comparez vos grilles : ce que l'un a coché « pas encore » devient son objectif prioritaire.\\
2. Modèle de réponse possible : \eng{Last month, I went to the immigration office in Yaoundé. First, I took a number and waited in the queue. Then, I filled in the form and paid the fee. After that, I brought my photos to the officer. Finally, I collected my passport two weeks later.} — Vérifiez : \eng{went}, \eng{took}, \eng{brought} (irréguliers) ; \eng{filled}, \eng{paid}, \eng{collected} (réguliers, attention à l'orthographe de \eng{paid}).\\
3. Parlez lentement, accentuez les syllabes fortes et terminez par une phrase conclusive : \eng{It was easier than I expected!} « C'était plus facile que je ne le pensais ! »
\end{corrige}

Vous savez désormais saluer, annoncer votre but, poser des questions, donner vos informations, vérifier votre compréhension, négocier une solution et conclure poliment — puis raconter toute la démarche au passé. Autant de compétences qui vous serviront bien au-delà de la salle de classe, au premier guichet venu.

\newpage

\coursec{Activité d'intégration}

\courssub{Objectifs de l'activité}

À l'issue de cette activité d'intégration, tu seras capable de :

\begin{itemize}
\item mobiliser dans un \cle{même échange oral} l'ensemble des savoir-faire de la leçon : saluer, demander une information, donner une information, relancer son interlocuteur et conclure poliment ;
\item comprendre et utiliser le \cle{vocabulaire administratif} lié à la demande de passeport (\eng{application form, birth certificate, receipt, processing time}...) ;
\item employer correctement les structures grammaticales étudiées : questions avec \eng{Wh-} et inversion, \eng{would like}, \eng{could you...?}, futur avec \eng{will} pour les délais ;
\item soigner ta \cle{prononciation}, ton accentuation et ton intonation (intonation montante pour les questions \eng{yes/no}, descendante pour les questions en \eng{Wh-}) ;
\item écouter activement le dialogue des autres binômes et relever des informations précises dans une grille d'écoute.
\end{itemize}

\courssub{Situation}

La classe est transformée en guichet de la \cle{Délégation Générale de la Sûreté Nationale} (DGSN) de Yaoundé. C'est la période des grandes vacances : de nombreux élèves et étudiants viennent demander un passeport pour voyager, poursuivre leurs études à l'étranger ou participer à des compétitions sportives. Au guichet, un agent d'accueil renseigne les usagers, vérifie leurs dossiers et leur indique les tarifs et les délais.

Tu vas jouer l'une des deux scènes suivantes avec ton camarade, à partir des fiches de rôle distribuées par l'enseignant.

\begin{textebox}[Role card A — The applicant]
You are \textbf{Mballa Sandrine}, a final-year student at Lycée de Nkol-Eton, Yaoundé. You were born on 12 March 2007 in Bafoussam. You have just won a scholarship to study in Ghana and you need a passport urgently. You have your birth certificate and your national ID card, but you have \textbf{no passport photos} and you do not know the price or the processing time. Ask the officer at least four questions. Be polite, greet the officer, and thank him or her at the end.
\end{textebox}

\begin{textebox}[Role card B — The passport officer]
You are the reception officer at the DGSN passport office, Yaoundé. Greet the applicant and answer his or her questions. Official information: required documents are a certified birth certificate, a national ID card, four passport photos and the application form filled in and signed. The fee is 75,000 FCFA for an ordinary passport. Normal processing time is two weeks; an express service (five working days) costs 110,000 FCFA. If a document is missing, explain clearly what the applicant must do and invite him or her to come back.
\end{textebox}

\begin{textebox}[Listening grid — for the audience]
While you listen to your classmates' dialogue, tick what you hear and complete the missing information.
\begin{itemize}
\item Documents mentioned: birth certificate / ID card / photos / form / receipt
\item Price of the ordinary passport: \ldots\ldots\ldots\ldots FCFA
\item Processing time (normal / express): \ldots\ldots\ldots\ldots / \ldots\ldots\ldots\ldots
\item Polite expressions heard: \eng{Good morning} / \eng{Could you...?} / \eng{I would like...} / \eng{Thank you very much}
\item Problem mentioned by the applicant: \ldots\ldots\ldots\ldots
\end{itemize}
\end{textebox}

\courssub{Consignes}

\begin{enumerate}
\item \textbf{Prépare ton rôle (5 minutes).} Lis attentivement ta fiche de rôle. Souligne les informations que tu dois \emph{donner} et celles que tu dois \emph{demander}. Note sur une feuille au moins quatre questions (demandeur) ou quatre informations clés (agent), sans rédiger le dialogue en entier : tu dois rester capable d'improviser.

\item \textbf{Révise les formules fonctionnelles (2 minutes).} Avec ton binôme, rappelez-vous une formule pour chaque fonction : saluer, demander une information, donner une information, relancer (\eng{Sorry, could you repeat that, please?}), conclure.

\item \textbf{Improvise le dialogue.} Joue la scène avec ton binôme, debout, face à face, comme au guichet. Ton dialogue doit comporter \cle{au moins dix répliques} (cinq par élève), contenir au moins quatre questions, mentionner les documents, le prix et le délai, et se terminer par une formule de politesse.

\item \textbf{Écoute les autres binômes.} Pendant que tes camarades jouent leur scène, remplis ta grille d'écoute : documents mentionnés, prix, délai, formules de politesse entendues, problème évoqué.

\item \textbf{Participe au feedback collectif.} Compare ta grille avec celle de tes voisins. Signale à la classe une formule bien employée et une erreur de langue ou de prononciation que tu as remarquée (sans nommer l'élève).

\item \textbf{Échange les rôles.} Recommence la scène en inversant les rôles, avec une nouvelle fiche de rôle (par exemple : un commerçant de Douala qui a perdu son reçu de paiement, ou un élève de Bamenda dont l'acte de naissance comporte une erreur de date).
\end{enumerate}

\courssub{Ressources à mobiliser}

\begin{aretenir}[Formules fonctionnelles à réutiliser]
\begin{itemize}
\item \textbf{Saluer :} \eng{Good morning, Sir/Madam. How can I help you?}
\item \textbf{Demander une information :} \eng{I would like to apply for a passport.} / \eng{What documents do I need?} / \eng{How much does it cost?} / \eng{How long does it take?} / \eng{Could you tell me where I can get the photos taken?}
\item \textbf{Donner une information :} \eng{You will need a certified birth certificate.} / \eng{The fee is 75,000 FCFA.} / \eng{It will be ready in two weeks.}
\item \textbf{Relancer :} \eng{Sorry, could you repeat that, please?} / \eng{Do you mean four photos?}
\item \textbf{Conclure :} \eng{Thank you very much for your help.} / \eng{You are welcome. Come back with your photos, please.}
\end{itemize}
\end{aretenir}

\begin{propriete}[Rappels grammaticaux utiles]
\begin{itemize}
\item Questions en \eng{Wh-} : mot interrogatif + auxiliaire + sujet + verbe. \eng{How long does it take?} (« Combien de temps cela prend-il ? ») — jamais \eng{How long it takes?} dans une question directe.
\item \eng{Would like} + nom ou infinitif pour une demande polie : \eng{I would like to apply for a passport.}
\item \eng{Could you} + base verbale : \eng{Could you repeat that, please?}
\item \eng{Will} pour annoncer un délai : \eng{Your passport will be ready in two weeks.}
\end{itemize}
\end{propriete}

\begin{attention}[Prononciation et erreurs à éviter]
\begin{itemize}
\item Intonation \cle{montante} pour les questions fermées : \eng{Do you have your ID card?} ; intonation \cle{descendante} pour les questions en \eng{Wh-} : \eng{What documents do I need?}
\item \eng{passport} se prononce « \textbf{pass}-porte » (accent sur la première syllabe) ; \eng{certificate} « seur-\textbf{ti}-fi-kète » ; \eng{receipt} « ri-\textbf{ssite} » (le \eng{p} ne se prononce pas).
\item Faux ami : \eng{to apply for} signifie « faire une demande de », pas « appliquer » au sens d'« appliquer une crème ».
\item Ne dis pas \eng{informations} au pluriel : \eng{information} est indénombrable.
\end{itemize}
\end{attention}

\courssub{Proposition de production et corrigé commenté}

\begin{exoresolu}[Dialogue modèle : At the Passport Office]
Joue ou lis le dialogue suivant, puis observe le commentaire des choix de langue.

\textbf{Officer:} Good morning, Madam. How can I help you?

\textbf{Sandrine:} Good morning, Sir. I would like to apply for a passport, please.

\textbf{Officer:} Certainly. Is it your first application?

\textbf{Sandrine:} Yes, it is. What documents do I need?

\textbf{Officer:} You will need a certified birth certificate, your national ID card, four passport photos and the application form, filled in and signed.

\textbf{Sandrine:} I have my birth certificate and my ID card, but I am afraid I do not have any photos. Could you tell me where I can get them taken?

\textbf{Officer:} There is a photo studio just opposite our office. It takes about ten minutes.

\textbf{Sandrine:} Thank you. And how much does an ordinary passport cost?

\textbf{Officer:} The fee is 75,000 FCFA, and the processing time is two weeks.

\textbf{Sandrine:} Two weeks? Sorry, could you repeat that, please?

\textbf{Officer:} Yes, two weeks. But if you are in a hurry, there is an express service: five working days, for 110,000 FCFA.

\textbf{Sandrine:} I see. I will think about it. Can I come back this afternoon with my photos?

\textbf{Officer:} Of course. We close at half past three.

\textbf{Sandrine:} Perfect. Thank you very much for your help, Sir.

\textbf{Officer:} You are welcome, Madam. See you later.

\tcblower

\textbf{Commentaire des choix de langue :}

\begin{itemize}
\item \textbf{Salutations et politesse :} le dialogue s'ouvre avec \eng{Good morning} et la question rituelle de l'agent \eng{How can I help you?} ; il se clôt avec \eng{Thank you very much} / \eng{You are welcome}. Les titres \eng{Sir} et \eng{Madam} marquent le respect, indispensable dans un contexte administratif camerounais.
\item \textbf{Demande polie :} \eng{I would like to apply for a passport} est plus courtois que \eng{I want a passport}. De même, \eng{Could you tell me where I can get them taken?} montre la question indirecte : après \eng{Could you tell me}, on revient à l'ordre sujet + verbe (\eng{where I can get}), sans inversion.
\item \textbf{Questions directes bien formées :} \eng{What documents do I need?} et \eng{How much does an ordinary passport cost?} respectent l'inversion avec l'auxiliaire \eng{do/does}.
\item \textbf{Annonce des délais :} \eng{You will need...}, \eng{It will be ready...} emploient \eng{will} pour une information certaine.
\item \textbf{Relance :} \eng{Sorry, could you repeat that, please?} permet de gérer un problème de compréhension sans bloquer l'échange — stratégie essentielle à l'oral.
\item \textbf{Gestion du problème :} \eng{I am afraid I do not have any photos} introduit poliment une mauvaise nouvelle ; l'agent propose alors une solution concrète, ce qui rend le dialogue réaliste.
\item \textbf{Prononciation :} veille à l'accentuation : « \textbf{pass}-porte », « \textbf{ser}-tifaid », « \textbf{ex}-press », et à l'intonation descendante sur \eng{What documents do I need?}
\end{itemize}
\end{exoresolu}

\courssub{Bilan de l'activité}

\begin{aretenir}[Ce que je sais faire maintenant]
À l'issue de cette activité, vérifie que tu peux :
\begin{itemize}
\item ouvrir et clore poliment un échange administratif en anglais ;
\item poser au moins quatre types de questions (\eng{What...?}, \eng{How much...?}, \eng{How long...?}, \eng{Could you...?}) avec l'inversion correcte ;
\item donner une information complète : documents, tarif, délai ;
\item relancer ton interlocuteur sans interrompre la communication ;
\item prononcer correctement le vocabulaire clé et placer l'intonation juste selon le type de question.
\end{itemize}
\end{aretenir}

\begin{methode}[Pour aller plus loin : la fiche d'auto-évaluation]
Après ton passage au « guichet », note-toi de 0 à 2 sur chaque critère : salutations et conclusion ; nombre et correction des questions ; exactitude des informations données ; utilisation d'au moins une relance ; prononciation et intonation ; aisance générale. Un total de 10 à 12 indique une maîtrise satisfaisante ; en dessous de 7, rejoue la scène avec un nouveau partenaire en t'appuyant sur la boîte à outils de la leçon.
\end{methode}

\begin{savaistu}[Le passeport camerounais]
Au Cameroun, les passeports sont délivrés par la Délégation Générale de la Sûreté Nationale (DGSN). Depuis 2021, le passeport biométrique est produit dans un centre unique à Yaoundé, et les demandes se font en ligne avant le retrait dans les antennes régionales, notamment à Douala, Bafoussam, Garoua ou Buea. Savoir échanger en anglais dans ce contexte est un atout réel : l'anglais est l'une des deux langues officielles du pays, et il est indispensable pour voyager dans les pays anglophones voisins comme le Nigeria ou le Ghana.
\end{savaistu}

\newpage

\coursec{Méthodes et exercices résolus}

\courssub{Fiches méthode et exercices résolus}

\begin{methode}[Échanger des informations dans la situation proposée]
Pour mener un échange d'informations sur une demande de passeport, suis toujours la même structure en cinq temps :
\begin{enumerate}
\item \textbf{Saluer et se présenter} : \eng{Good morning, sir. I would like to apply for a passport.} — commence toujours par une formule de politesse (\eng{Good morning}, \eng{Excuse me}).
\item \textbf{Annoncer clairement ton objectif} : utilise \eng{I would like to + base verbale} ou \eng{I need some information about + nom}. Évite le calque \emph{I want} (trop direct, impoli dans un contexte administratif).
\item \textbf{Poser des questions précises} : une question à la fois, avec inversion sujet–auxiliaire : \eng{What documents do I need?} \eng{How long does it take?} \eng{How much does it cost?}
\item \textbf{Réagir et reformuler} : montre que tu écoutes (\eng{I see.}, \eng{All right.}) et reformule pour vérifier : \eng{So I need my birth certificate and two photos, is that right?}
\item \textbf{Conclure poliment} : \eng{Thank you very much for your help. Goodbye.}
\end{enumerate}
\textbf{Réflexes à retenir} : \eng{apply for} (jamais \emph{apply a passport}), \eng{fill in a form}, \eng{hand in / submit an application}, \eng{How long does it take?} (jamais \emph{How much time it takes?}).
\end{methode}

\begin{exoresolu}[Échanger des informations — dialogue à compléter]
\textbf{Énoncé.} Complete the dialogue between a student and an immigration officer in Yaoundé with appropriate sentences.\par
\eng{A: Good morning, madam. (1) ....................................................?}\par
\eng{B: You need your birth certificate, a national ID card and two passport photos.}\par
\eng{A: (2) ....................................................?}\par
\eng{B: It costs 75,000 francs CFA.}\par
\eng{A: (3) ....................................................?}\par
\eng{B: It usually takes two weeks.}\par
\eng{A: Thank you very much for your help. Goodbye.}
\tcblower
\textbf{Raisonnement.} Chaque réponse de l'officier (B) révèle la nature de la question posée par l'élève (A) : il faut donc partir des réponses pour reconstruire les questions, avec le bon mot interrogatif et l'inversion sujet–auxiliaire.\par
(1) La réponse énumère des \emph{documents} → mot interrogatif \eng{What} + auxiliaire \eng{do} : \eng{What documents do I need to apply for a passport?}\par
(2) La réponse donne un \emph{prix} → \eng{How much} + inversion : \eng{How much does it cost?} (attention : \eng{cost} reste à la base verbale car \eng{does} porte déjà la marque de la 3\textsuperscript{e} personne).\par
(3) La réponse donne une \emph{durée} → \eng{How long} + inversion : \eng{How long does it take to get the passport?}\par
\textbf{Conclusion.} Dans un dialogue à trous, lis d'abord toutes les réponses : elles dictent le mot interrogatif (\eng{what / how much / how long / where / when}) et le temps verbal de la question.
\end{exoresolu}

\begin{methode}[Poser des questions et répondre de façon adaptée]
\begin{enumerate}
\item \textbf{Choisir le bon mot interrogatif} : \eng{What} (chose), \eng{Where} (lieu), \eng{When} (moment), \eng{How much} (prix, quantité indénombrable), \eng{How many} (quantité dénombrable), \eng{How long} (durée), \eng{Why} (raison).
\item \textbf{Appliquer l'inversion} : mot interrogatif + auxiliaire (\eng{do/does/did}, \eng{can}, \eng{is/are}) + sujet + verbe. \eng{Where can I get the form?}
\item \textbf{Questions fermées (oui/non)} : auxiliaire + sujet + verbe, avec intonation montante : \eng{Do you accept photocopies?}
\item \textbf{Répondre de façon adaptée} : réponse courte + information complète. À \eng{Do I need an appointment?} réponds \eng{Yes, you do. You must book one online.}
\item \textbf{Relancer poliment} si la réponse est incomplète : \eng{Could you tell me more about the fees?} / \eng{Sorry, could you repeat that, please?}
\end{enumerate}
\textbf{Structure à retenir} : dans les questions indirectes, \textbf{pas d'inversion} : \eng{Could you tell me where the office is?} (jamais \emph{where is the office} après \eng{Could you tell me}).
\end{methode}

\begin{exoresolu}[Poser des questions — transformation]
\textbf{Énoncé.} Ask the questions that correspond to these answers.\par
1. \eng{The passport office opens at 7.30 a.m.} (When)\par
2. \eng{You need four passport photos.} (How many)\par
3. \eng{You can pay at the bank or by mobile money.} (Where)\par
4. \eng{I want a passport because I am going to study in Ghana.} (Why)\par
5. Put into an indirect question: \eng{Where is the application form?} → \eng{Could you tell me ...?}
\tcblower
\textbf{Solutions et justifications.}\par
1. \eng{When does the passport office open?} — présent simple (habitude/horaire), auxiliaire \eng{does} car le sujet est à la 3\textsuperscript{e} personne ; le verbe \eng{open} reste à la base verbale.\par
2. \eng{How many passport photos do I need?} — \eng{photos} est dénombrable pluriel, donc \eng{How many} (et non \emph{How much}).\par
3. \eng{Where can I pay?} — inversion du modal \eng{can} devant le sujet \eng{I}.\par
4. \eng{Why do you want a passport?} — \eng{Why} + \eng{do} + sujet + base verbale.\par
5. \eng{Could you tell me where the application form is?} — question indirecte : on revient à l'ordre sujet + verbe (\eng{the application form is}), sans inversion. C'est la règle la plus souvent sanctionnée au Bac.\par
\textbf{Conclusion.} Question directe = inversion ; question indirecte après \eng{Could you tell me / Do you know} = ordre normal sujet–verbe.
\end{exoresolu}

\begin{methode}[Mener un court dialogue en anglais — le jeu de rôle]
\begin{enumerate}
\item \textbf{Préparer les rôles} : identifie qui tu es (applicant / officer) et ce que tu dois obtenir (documents requis, prix, délai, lieu).
\item \textbf{Ouvrir l'échange} : salutation + objectif. \eng{Good morning. I would like to apply for a passport.}
\item \textbf{Alterner questions et réponses} : au moins trois échanges ; chaque réplique doit faire avancer la conversation (question → réponse → relance).
\item \textbf{Utiliser les formules de relance} : \eng{Could you repeat that, please?}, \eng{What do you mean by “supporting documents”?}, \eng{Is there anything else I should bring?}
\item \textbf{Conclure} : remercier + prendre congé : \eng{Thank you for your help. Have a nice day.}
\item \textbf{Soigner la prononciation} : accentue le mot porteur de sens (\eng{PASSport}, \eng{DOCuments}, \eng{appliCAtion} — « a-pli-ké-chion »), voix montante pour les questions fermées, descendante pour les questions en \eng{wh-}.
\end{enumerate}
\textbf{Réflexe} : parle en phrases complètes, jamais en mots isolés ; si tu bloques, gagne du temps avec \eng{Well...}, \eng{Let me see...}.
\end{methode}

\begin{exoresolu}[Mener un dialogue — rédaction d'un jeu de rôle]
\textbf{Énoncé.} You are at the immigration office in Douala. Write a short dialogue (at least six exchanges) in which you ask for information and then hand in your passport application.
\tcblower
\textbf{Démarche.} On suit la structure en cinq temps : salutation → objectif → questions (documents, prix, délai) → dépôt du dossier → conclusion. On varie les mots interrogatifs et on utilise \eng{would like}, \eng{need}, \eng{fill in}, \eng{hand in}.\par
\textbf{Réponse rédigée.}\par
\eng{A: Good morning, sir. I would like to apply for a passport.}\par
\eng{B: Good morning. Have you filled in the application form?}\par
\eng{A: Yes, I have. Here it is. What other documents do you need?}\par
\eng{B: I need your birth certificate, your national ID card and two recent passport photos.}\par
\eng{A: Here they are. How much is the fee, please?}\par
\eng{B: It is 75,000 francs CFA. You can pay at the counter over there.}\par
\eng{A: All right. And how long will it take to get the passport?}\par
\eng{B: About two weeks. We will send you a text message when it is ready.}\par
\eng{A: Perfect. Here is my application file.}\par
\eng{B: Thank you. Everything is complete.}\par
\eng{A: Thank you very much for your help. Goodbye, sir.}\par
\textbf{Justifications.} \eng{Have you filled in...?} : present perfect car l'action est achevée avec un résultat présent. \eng{How long will it take...?} : futur avec \eng{will} car on parle d'un délai à venir. Attention à \eng{Here it is / Here they are} : après \eng{here}, il n'y a \textbf{pas} d'inversion quand le sujet est un \emph{pronom} (\eng{Here it is}), mais il y a inversion avec un \emph{nom} (\eng{Here is your receipt}).\par
\textbf{Conclusion.} Un bon dialogue d'examen contient : une ouverture polie, au moins trois questions variées, des réactions (\eng{All right}, \eng{Perfect}) et une clôture.
\end{exoresolu}

\begin{methode}[Réagir aux questions de compréhension sur un dialogue]
\begin{enumerate}
\item \textbf{Lire les questions avant le texte} : repère les mots interrogatifs pour savoir quoi chercher (lieu, prix, durée, documents).
\item \textbf{Repérer les mots-clés} dans le dialogue : \eng{fee} = prix, \eng{take two weeks} = délai, \eng{supporting documents} = pièces justificatives.
\item \textbf{Répondre en phrases complètes} en reprenant les termes de la question : à \eng{How much does a passport cost?} réponds \eng{A passport costs 75,000 francs CFA.} — jamais un mot seul.
\item \textbf{Ne pas recopier tout le texte} : extrais uniquement l'information demandée.
\item \textbf{Vérifier la concordance des temps} : si la question est au présent, réponds au présent.
\end{enumerate}
\end{methode}

\begin{exoresolu}[Compréhension + vocabulaire]
\textbf{Énoncé.} Read the dialogue and answer the questions in complete sentences.\par
\eng{Officer: Good morning. Can I help you?}\par
\eng{Applicant: Yes, please. I would like to renew my passport. It expired last month.}\par
\eng{Officer: You need to fill in this form and bring your old passport, two photos and a photocopy of your ID card.}\par
\eng{Applicant: How much is the renewal fee?}\par
\eng{Officer: It is 75,000 francs, and the passport will be ready in ten working days.}\par
1. What does the applicant want to do?\par
2. How many photos must he bring?\par
3. When will the new passport be ready?\par
4. Find in the text a word that means “no longer valid”.
\tcblower
\textbf{Solutions.}\par
1. \eng{The applicant wants to renew his passport.} — on reprend \eng{renew} du texte ; \eng{want} prend un \eng{s} (3\textsuperscript{e} personne du présent simple).\par
2. \eng{He must bring two photos.} — information extraite de la liste des documents ; \eng{must} + base verbale.\par
3. \eng{The new passport will be ready in ten working days.} — futur \eng{will be ready} comme dans le texte ; \eng{working days} = jours ouvrables.\par
4. \eng{Expired.} — \eng{It expired last month} signifie « il a expiré le mois dernier » ; \eng{expire} = « ne plus être valide » (à ne pas confondre avec le faux ami \emph{expirer} au sens de « mourir »).\par
\textbf{Conclusion.} Chaque réponse reprend la structure de la question et un seul élément du texte : précision et phrases complètes sont les deux critères de notation.
\end{exoresolu}

\begin{methode}[Soigner prononciation, accentuation et intonation à l'oral]
\begin{enumerate}
\item \textbf{Accentuer correctement les mots clés} : \eng{PASSport} (« pâss-port »), \eng{appliCAtion} (« a-pli-ké-chion »), \eng{DOCument} (« do-kieu-ment »), \eng{cerTIficate} (« seur-ti-fi-ket »), \eng{PHOtograph} (« fo-to-graf »).
\item \textbf{Intonation montante} (la voix monte en fin de phrase) pour les questions fermées : \eng{Do you accept photocopies? ↗}
\item \textbf{Intonation descendante} pour les questions en \eng{wh-} et les affirmations : \eng{What documents do I need? ↘}
\item \textbf{Prononcer les finales} : le \eng{-s} de la 3\textsuperscript{e} personne (\eng{it costs} — « its kosts ») et le \eng{-ed} du prétérit (\eng{expired} — « èk-spaïeurd »).
\item \textbf{Attention au « th »} : \eng{thank you} (« sènk iou », langue entre les dents), jamais « zank you ».
\end{enumerate}
\end{methode}

\begin{exoresolu}[Prononciation et intonation]
\textbf{Énoncé.} 1. Mark the stressed syllable in: \eng{passport, application, certificate, photograph, information}. 2. Say whether the intonation rises (↗) or falls (↘): (a) \eng{Where is the passport office?} (b) \eng{Do you have an appointment?} (c) \eng{Your passport will be ready soon.}
\tcblower
\textbf{Solutions.}\par
1. \eng{PASSport} ; \eng{appliCAtion} ; \eng{cerTIficate} ; \eng{PHOtograph} ; \eng{inforMAtion}. Règle utile : les mots en \eng{-tion} portent l'accent sur la syllabe qui précède (\eng{appliCAtion}, \eng{inforMAtion} — « a-pli-ké-chion », « in-for-mé-chion »).\par
2. (a) ↘ intonation descendante : question ouverte en \eng{wh-}. (b) ↗ intonation montante : question fermée appelant \eng{yes/no}. (c) ↘ intonation descendante : phrase affirmative, la voix descend en fin de phrase.\par
\textbf{Conclusion.} À l'oral du Bac, une intonation correcte montre la maîtrise de la communication : monte pour les questions oui/non, descends pour les questions en \eng{wh-} et les réponses.
\end{exoresolu}

\begin{exercice}[À toi de jouer — dialogue et questions]
1. Complete with the correct question word (\eng{what, where, when, how much, how many, how long, why}):\par
(a) \eng{.............. does it take to get a passport? — About two weeks.}\par
(b) \eng{.............. photos do you need? — Two recent ones.}\par
(c) \eng{.............. can I collect my passport? — At this office.}\par
2. Put the words in the correct order to make polite questions:\par
(a) \eng{is / the / where / office / passport / the / ?}\par
(b) \eng{tell / could / me / you / costs / how much / it / ?}\par
3. With a partner, act out the situation: one of you is an applicant at the immigration office in Bamenda, the other is the officer. Exchange at least four questions and answers about documents, fees and delay.
\end{exercice}

\begin{corrige}[Exercice : À toi de jouer]
1. (a) \eng{How long} (durée) ; (b) \eng{How many} (dénombrable pluriel) ; (c) \eng{Where} (lieu).\par
2. (a) \eng{Where is the passport office?} — question directe avec le verbe \eng{be} : inversion verbe–sujet, sans auxiliaire \eng{do}. (b) \eng{Could you tell me how much it costs?} — question indirecte : ordre sujet–verbe (\eng{it costs}), avec le \eng{-s} de la 3\textsuperscript{e} personne conservé.\par
3. Exemple de production attendue : salutation polie, \eng{I would like to apply for a passport}, questions variées (\eng{What documents...?}, \eng{How much...?}, \eng{How long...?}), réactions (\eng{I see}, \eng{All right}) et conclusion (\eng{Thank you for your help}).
\end{corrige}

\begin{attention}[Les 5 erreurs qui coûtent des points au Bac]
\begin{enumerate}
\item \textbf{Calque de la question française sans inversion} : \emph{How much it costs?} est faux. Écris \eng{How much does it cost?} — auxiliaire \eng{does} + base verbale, sans \eng{s} sur \eng{cost}.
\item \textbf{Inversion dans la question indirecte} : \emph{Could you tell me where is the office?} est faux. Après \eng{Could you tell me / Do you know}, ordre sujet–verbe : \eng{Could you tell me where the office is?}
\item \textbf{Faux amis et mauvaise préposition} : on dit \eng{apply \textbf{for} a passport} (jamais \emph{apply a passport}), \eng{fill \textbf{in} a form} (et non \emph{fill a form}), \eng{wait \textbf{for} the passport}. Et \emph{actually} ne veut pas dire « actuellement » : dis \eng{currently}.
\item \textbf{Confusion how much / how many et oubli du \eng{-s}} : \eng{How much does it cost?} (prix) mais \eng{How many photos do I need?} (dénombrable). N'oublie pas le \eng{-s} de la 3\textsuperscript{e} personne : \eng{It cost\textbf{s} 75,000 francs.}
\item \textbf{Orthographe et prononciation des mots clés} : \eng{passport} (un seul \eng{p} au milieu : \emph{passport}, pas \emph{passeport}), \eng{certificate} (pas \emph{certificat}), \eng{appointment} (deux \eng{p}), \eng{government} (avec le \eng{n} : « go-veurn-ment »). Une faute sur ces mots du thème est toujours sanctionnée.
\end{enumerate}
\end{attention}

\newpage

\coursec{Exercices}

\courssub{Je vérifie mes connaissances}

\begin{exercice}[QCM : les bons réflexes au guichet]
Pour chaque question, choisis la bonne réponse (a, b ou c).
\begin{enumerate}
\item To ask politely for a form, you say :
\begin{enumerate}
\item “Give me the form now.”
\item “Could I have the form, please?”
\item “I want form.”
\end{enumerate}
\item “How long does it take to process the application?” means :
\begin{enumerate}
\item Combien coûte la demande ?
\item Combien de temps prend le traitement de la demande ?
\item Où doit-on déposer la demande ?
\end{enumerate}
\item To thank the officer at the end of the exchange, you say :
\begin{enumerate}
\item “You're welcome.”
\item “Thank you very much for your help.”
\item “Excuse me, please.”
\end{enumerate}
\item A “receipt” is :
\begin{enumerate}
\item un récépissé / une quittance de paiement
\item un formulaire de demande
\item une photo d'identité
\end{enumerate}
\end{enumerate}
\end{exercice}

\begin{exercice}[Vrai ou faux — justifie ta réponse]
Réponds par \emph{True} ou \emph{False}, puis justifie en une phrase en anglais.
\begin{enumerate}
\item “Could you tell me where the immigration office is?” is a polite way to ask for information.
\item In English, you can say “I am agree with this fee.”
\item “How much is the passport fee?” is a correct question about the price.
\item When the officer says “Thank you”, the applicant should answer “Please”.
\end{enumerate}
\end{exercice}

\begin{exercice}[Vocabulaire : le dossier de demande]
Complète chaque phrase avec le mot anglais qui convient, choisi dans la liste : \emph{application form — birth certificate — passport photos — receipt — fee — embassy}.
\begin{enumerate}
\item You must fill in an \ldots\ldots\ldots with your name, date of birth and address.
\item Two recent \ldots\ldots\ldots are required for the file.
\item The officer will give you a \ldots\ldots\ldots after you pay.
\item The passport \ldots\ldots\ldots is paid at the cashier's desk.
\item A certified copy of your \ldots\ldots\ldots proves your date and place of birth.
\item If you are abroad, you apply at your country's \ldots\ldots\ldots.
\end{enumerate}
\end{exercice}

\begin{exercice}[Conjugaison et formes à compléter]
Complète les phrases avec la forme correcte du verbe entre parenthèses.
\begin{enumerate}
\item I \ldots\ldots\ldots (to want) to apply for a passport, please.
\item How long \ldots\ldots\ldots it \ldots\ldots\ldots (to take) to get the passport?
\item You \ldots\ldots\ldots (must / to fill) in this form in capital letters.
\item She \ldots\ldots\ldots (to bring) all her documents yesterday.
\item The passports \ldots\ldots\ldots (to be) ready in two weeks.
\item \ldots\ldots\ldots you \ldots\ldots\ldots (can / to tell) me where the photo studio is?
\end{enumerate}
\end{exercice}

\courssub{Je m'entraîne}

\begin{exercice}[Traduction : au service de l'immigration]
Traduis les phrases suivantes en anglais.
\begin{enumerate}
\item Bonjour Monsieur, je voudrais faire une demande de passeport.
\item Combien de photos d'identité faut-il fournir ?
\item Vous devez payer les frais à la caisse et revenir avec le reçu.
\item Quand est-ce que je peux retirer mon passeport ?
\item Excusez-moi, pourriez-vous répéter, s'il vous plaît ?
\end{enumerate}
\end{exercice}

\begin{attention}[Faux amis et pièges de traduction]
\begin{itemize}
\item « Faire une demande de passeport » se dit \eng{to apply for a passport} (jamais \eng{to make a demand} : \emph{demand} en anglais signifie « exigence »).
\item « Retirer » un document se dit \eng{to collect} (ou \eng{to pick up}), pas \eng{to remove}.
\item « Fournir » se dit \eng{to provide} ; attention, \eng{to furnish} signifie « meubler ».
\item « Combien de » + nom dénombrable = \eng{How many} (photos) ; + nom indénombrable ou prix = \eng{How much} (money, the fee).
\end{itemize}
\end{attention}

\begin{exercice}[Des demandes directes aux demandes polies]
Transforme chaque demande directe en demande polie, en utilisant \emph{Could you\ldots?}, \emph{Could I\ldots?} ou \emph{Would you mind\ldots?}.
\begin{enumerate}
\item Give me the form.
\item Tell me the fee.
\item Repeat your name.
\item Show me the list of documents.
\item Wait for me outside.
\end{enumerate}
\end{exercice}

\begin{exercice}[Dialogue à trous : déposer sa demande]
Complète le dialogue ci-dessous avec les formules fonctionnelles de la liste : \emph{How long does it take? — Could I have an application form, please? — You're welcome. — Here you are. — Thank you very much for your help. — How much is the fee?}
\end{exercice}

\begin{textebox}[At the immigration office]
Applicant: Good morning, Madam. \ldots\ldots\ldots\ldots\ldots\ldots\ldots\ldots

Officer: Certainly. \ldots\ldots\ldots\ldots\ldots\ldots\ldots\ldots Please fill it in carefully.

Applicant: \ldots\ldots\ldots\ldots\ldots\ldots\ldots\ldots

Officer: It is 75,000 FCFA, payable at the cashier's desk.

Applicant: And \ldots\ldots\ldots\ldots\ldots\ldots\ldots\ldots

Officer: About two weeks. Come back with your receipt.

Applicant: \ldots\ldots\ldots\ldots\ldots\ldots\ldots\ldots

Officer: \ldots\ldots\ldots\ldots\ldots\ldots\ldots\ldots Have a nice day.
\end{textebox}

\begin{exercice}[Appariement : questions et réponses]
Associe chaque question du demandeur (1--6) à la réponse de l'agent (a--f).
\begin{enumerate}
\item What documents do I need?
\item How much is the passport fee?
\item How long does the process take?
\item Where can I pay?
\item When can I collect my passport?
\item Could you repeat that, please?
\end{enumerate}
\begin{itemize}
\item[a)] At the cashier's desk, on the ground floor.
\item[b)] Two weeks from today, Madam.
\item[c)] Of course: I said you need two passport photos.
\item[d)] An application form, a birth certificate and two photos.
\item[e)] It costs 75,000 FCFA.
\item[f)] Next Friday, with your receipt.
\end{itemize}
\end{exercice}

\courssub{Je me prépare au Bac}

\begin{exercice}[Compréhension écrite et langue — type Bac]
Lis le texte suivant, puis réponds aux questions en anglais, par des phrases complètes.
\end{exercice}

\begin{textebox}[A morning at the passport office in Yaoundé]
Last Tuesday, Amina Mbarga, a final-year student from Lycée de Nkol-Eton, went to the immigration office in Yaoundé to apply for her first passport. She had prepared her file carefully: a completed application form, a certified copy of her birth certificate, two recent passport photos and her national identity card.

At the reception desk, she greeted the officer politely and asked: “Could you tell me if my file is complete, please?” The officer checked every document and smiled. “Everything is in order,” he said. “You only need to pay the fee at the cashier's desk and bring me the receipt.”

Amina paid the fee and returned with the receipt. “How long will it take?” she asked. “About three weeks,” the officer replied. “We will send you a text message when your passport is ready. Come back with this receipt and your identity card.” Amina thanked him warmly and left the office, already dreaming of her first trip abroad.
\end{textebox}

\begin{exercice}[Questions sur le texte]
\textbf{A. Comprehension questions}
\begin{enumerate}
\item Who is Amina Mbarga and where does she study?
\item List the four documents Amina brought in her file.
\item What did the officer ask Amina to do after checking her documents?
\item How will Amina know that her passport is ready?
\item What must she bring when she comes back to collect the passport?
\end{enumerate}

\textbf{B. Language work}
\begin{enumerate}
\item Find in the text two polite expressions used by Amina.
\item Turn into a polite question: “Tell me if my file is complete.”
\item Give the French equivalent of: “Everything is in order.”
\end{enumerate}
\end{exercice}

\begin{exercice}[Traduction — type Bac]
Traduis le passage suivant en anglais.
\end{exercice}

\begin{textebox}[Texte à traduire]
Hier matin, mon frère est allé au bureau de l'immigration pour renouveler son passeport expiré. L'agent lui a demandé de remplir un formulaire et de fournir deux photos d'identité récentes. « Combien de temps cela prendra-t-il ? » a demandé mon frère. L'agent a répondu qu'il pourrait retirer son nouveau passeport dans trois semaines, à condition de présenter son reçu.
\end{textebox}

\begin{exercice}[Expression écrite et orale — type Bac]
\textbf{Sujet.} Your cousin, who lives in Bamenda, wants to apply for a passport for the first time and has asked you for advice.

\textbf{Tâche écrite.} Write a short dialogue (at least 8 turns) between an applicant and an immigration officer at the passport office. Use at least four functional expressions studied in the lesson (asking for information, asking about the fee, asking about the delay, thanking and concluding). (80--100 words.)

\textbf{Tâche orale.} Be ready to play your dialogue with a classmate, and to answer the examiner's personal questions in complete sentences, as if you were at the desk: \emph{What is your name? When and where were you born? What is your occupation? What is your address?}
\end{exercice}



\newpage

\coursec{Corrigés des exercices}

\courssub{Corrigés de la partie « Je vérifie mes connaissances »}

\begin{corrige}[Exercice 1 — QCM : les bons réflexes au guichet]
\begin{enumerate}
\item \textbf{b)} \eng{Could I have the form, please?} — c'est la formule polie de demande avec \eng{Could I...?}. La réponse a) est un ordre brutal, inacceptable au guichet ; la réponse c) est fautive : il manque l'article indéfini (\eng{I want \cle{a} form} serait grammaticalement correct mais resterait trop direct).
\item \textbf{b)} \eng{How long does it take...?} interroge sur la \emph{durée} (« combien de temps »). Le prix se demanderait avec \eng{How much is the application?} et le lieu avec \eng{Where...?}.
\item \textbf{b)} \eng{Thank you very much for your help.} — attention : \eng{You're welcome} (« je vous en prie ») sert à \emph{répondre} à un remerciement, jamais à remercier ; \eng{Excuse me} sert à attirer l'attention ou à s'excuser.
\item \textbf{a)} un récépissé / une quittance de paiement. Le formulaire de demande est \eng{an application form} et la photo d'identité \eng{a passport photo}.
\end{enumerate}
\end{corrige}

\begin{corrige}[Exercice 2 — Vrai ou faux — justifie ta réponse]
\begin{enumerate}
\item \emph{True.} \eng{“Could you tell me where the immigration office is?” is a polite indirect question.} — la question indirecte introduite par \eng{Could you tell me...?} adoucit la demande ; attention à l'ordre des mots : sujet avant verbe après \eng{where} (\eng{where the office is}, pas \eng{where is the office}).
\item \emph{False.} \eng{“Agree” is a verb, so we say “I agree with this fee”, without “am”.} — \eng{to agree} est un verbe lexical : il ne prend pas l'auxiliaire \eng{be}. Erreur typique du francophone qui calque « je suis d'accord ».
\item \emph{True.} \eng{“How much...?” is used to ask about a price.} — \eng{How much is the passport fee?} est la question correcte sur le montant des frais.
\item \emph{False.} \eng{The applicant should answer “You're welcome”, not “Please”.} — à \eng{Thank you}, on répond \eng{You're welcome} ou \eng{My pleasure}.
\end{enumerate}
\end{corrige}

\begin{corrige}[Exercice 3 — Vocabulaire : le dossier de demande]
\begin{enumerate}
\item \textbf{\cle{application form}} — \eng{You must fill in an application form with your name, date of birth and address.} (« Vous devez remplir un formulaire de demande... »)
\item \textbf{\cle{passport photos}} — \eng{Two recent passport photos are required for the file.} (« Deux photos d'identité récentes sont exigées. »)
\item \textbf{\cle{receipt}} — \eng{The officer will give you a receipt after you pay.} (« L'agent vous remettra un récépissé après le paiement. »)
\item \textbf{\cle{fee}} — \eng{The passport fee is paid at the cashier's desk.} (« Les frais de passeport se paient à la caisse. »)
\item \textbf{\cle{birth certificate}} — \eng{A certified copy of your birth certificate proves your date and place of birth.} (« Une copie certifiée de votre acte de naissance... »)
\item \textbf{\cle{embassy}} — \eng{If you are abroad, you apply at your country's embassy.} (« À l'étranger, on fait sa demande auprès de l'ambassade de son pays. »)
\end{enumerate}
\end{corrige}

\begin{corrige}[Exercice 4 — Conjugaison et formes à compléter]
\begin{enumerate}
\item \textbf{want} — \eng{I want to apply for a passport, please.} Présent simple, première personne : base verbale. Registre plus poli : \eng{I \cle{would like} to apply for a passport, please.}
\item \textbf{does} it \textbf{take} — \eng{How long does it take to get the passport?} Question au présent simple : auxiliaire \eng{does} (3\textsuperscript{e} personne) + sujet + base verbale ; le verbe ne prend jamais de \eng{-s} après \eng{does}.
\item \textbf{must fill} — \eng{You must fill in this form in capital letters.} Après un modal (\eng{must}), base verbale sans \eng{to}.
\item \textbf{brought} — \eng{She brought all her documents yesterday.} Prétérit du verbe irrégulier \eng{to bring} (\emph{bring -- brought -- brought}), imposé par le marqueur \eng{yesterday}.
\item \textbf{will be} — \eng{The passports will be ready in two weeks.} Futur simple (\eng{will} + base verbale) pour une annonce, marqué par \eng{in two weeks}.
\item \textbf{Can} you \textbf{tell} — \eng{Can you tell me where the photo studio is?} Inversion du modal \eng{can} en tête de question + base verbale. Forme plus polie : \eng{Could you tell me...?}
\end{enumerate}
\end{corrige}

\courssub{Corrigés de la partie « Je m'entraîne »}

\begin{corrige}[Exercice 5 — Traduction : au service de l'immigration]
\begin{enumerate}
\item \eng{Good morning, Sir. I \cle{would like to apply for} a passport.} — « faire une demande de » = \eng{to apply for} ; éviter le calque \eng{to make a demand}.
\item \eng{How many passport photos must I provide?} (variante acceptée : \eng{How many passport photos are required?}) — \eng{photos} est dénombrable, donc \eng{How many}, pas \eng{How much}.
\item \eng{You must pay the fee at the \cle{cashier's desk} and come back with the receipt.} — « la caisse » = \eng{the cashier's desk} ; « revenir » = \eng{to come back}.
\item \eng{When can I \cle{collect} my passport?} — « retirer » un document = \eng{to collect} (ou \eng{to pick up}), jamais \eng{to remove}.
\item \eng{Excuse me, could you repeat that, please?} — \eng{Could you...?} pour la politesse ; \eng{repeat} se construit sans \eng{again} (\eng{repeat again} est un pléonasme à éviter).
\end{enumerate}
\end{corrige}

\begin{corrige}[Exercice 6 — Des demandes directes aux demandes polies]
\begin{enumerate}
\item \eng{Could I have the form, please?} (ou \eng{Could you give me the form, please?})
\item \eng{Could you tell me the fee, please?}
\item \eng{Could you repeat your name, please?}
\item \eng{Could you show me the list of documents, please?}
\item \eng{Would you mind waiting for me outside, please?} — point de grammaire : après \eng{Would you mind}, le verbe est obligatoirement à la forme en \emph{-ing} (\eng{waiting}, pas \eng{to wait}).
\end{enumerate}
Dans tous les cas, le \eng{please} final et l'intonation montante renforcent la politesse.
\end{corrige}

\begin{corrige}[Exercice 7 — Dialogue à trous : déposer sa demande]
\begin{enumerate}
\item \eng{Could I have an application form, please?} — le demandeur ouvre l'échange par une demande polie.
\item \eng{Here you are.} — formule standard de l'agent qui tend le document (« voici »).
\item \eng{How much is the fee?} — la réponse chiffrée (\eng{It is 75,000 FCFA}) impose une question sur le prix.
\item \eng{how long does it take?} — sans majuscule, car la question est introduite par \eng{And} ; la réponse \eng{About two weeks} confirme qu'il s'agit de la durée.
\item \eng{Thank you very much for your help.} — le demandeur conclut en remerciant.
\item \eng{You're welcome.} — réponse rituelle de l'agent au remerciement.
\end{enumerate}
\end{corrige}

\begin{corrige}[Exercice 8 — Appariement : questions et réponses]
\begin{center}
\begin{tabularx}{\linewidth}{|X|X|X|}
\hline
\rowcolor{popblueL} Question & Réponse & Justification \\
\hline
1. What documents do I need? & d & la réponse énumère des documents \\
\hline
2. How much is the passport fee? & e & \eng{How much} appelle un montant \\
\hline
3. How long does the process take? & b & \eng{How long} appelle une durée \\
\hline
4. Where can I pay? & a & \eng{Where} appelle un lieu \\
\hline
5. When can I collect my passport? & f & \eng{When} appelle une date \\
\hline
6. Could you repeat that, please? & c & l'agent répète sa consigne \\
\hline
\end{tabularx}
\end{center}
Réponses : 1 -- d ; 2 -- e ; 3 -- b ; 4 -- a ; 5 -- f ; 6 -- c.
\end{corrige}

\courssub{Corrigés de la partie « Je me prépare au Bac »}

\begin{corrige}[Exercice 9 — Compréhension écrite et langue — type Bac]
\textbf{A. Comprehension questions} (barème indicatif : 2 points par réponse correcte, en phrase complète)
\begin{enumerate}
\item \eng{Amina Mbarga is a final-year student, and she studies at Lycée de Nkol-Eton.}
\item \eng{She brought a completed application form, a certified copy of her birth certificate, two recent passport photos and her national identity card.}
\item \eng{The officer asked her to pay the fee at the cashier's desk and to bring him the receipt.}
\item \eng{She will know that her passport is ready because the office will send her a text message.}
\item \eng{When she comes back, she must bring her receipt and her identity card.}
\end{enumerate}
\textbf{B. Language work} (barème indicatif : 2 points par item)
\begin{enumerate}
\item Deux marques de politesse (énoncées par les personnages) : \eng{“Could you tell me if my file is complete, please?”} et \eng{“Thank you very much for your help.”} (ou \eng{“Good morning, Sir.”}).
\item \eng{Could you tell me if my file is complete, please?} — attention à l'ordre des mots de la question indirecte : \eng{if my file is complete} (sujet + verbe).
\item Anglais : \eng{“Your file is complete.”} / \eng{“Everything is in order.”} — Français : « Tout est en règle. » / « Votre dossier est complet. »
\end{enumerate}
Points de vigilance : répondre par des phrases complètes reprenant les termes de la question ; ne pas recopier tout le texte ; respecter la concordance des temps (le récit est au prétérit).
\end{corrige}

\begin{corrige}[Exercice 10 — Traduction — type Bac]
Proposition de traduction (barème indicatif : 10 points — 2 points par segment ; 1 point si le sens est correct malgré une faute mineure) :

\eng{Yesterday morning, my brother went to the immigration office to renew his expired passport. The officer asked him to fill in a form and to provide two recent passport photos. “How long will it take?” my brother asked. The officer replied that he could collect his new passport in three weeks, provided that he showed his receipt.}

Points de vigilance :
\begin{itemize}
\item Prétérit des verbes irréguliers : \eng{go -- went}, \eng{take -- took} ; « est allé » = \eng{went} (et non \eng{goed}).
\item « Renouveler » = \eng{to renew} ; « expiré » = \eng{expired}.
\item « Remplir un formulaire » = \eng{to \cle{fill in} a form} ; « fournir » = \eng{to \cle{provide}} (et non \eng{to furnish}).
\item Discours rapporté au passé : \eng{will} recule en \cle{would/could} — \eng{he could collect his new passport in three weeks}.
\item « À condition de » = \cle{provided that} + proposition (ou \cle{as long as}) ; concordance des temps : \eng{showed} (prétérit modal / \emph{backshift}).
\end{itemize}
\end{corrige}

\begin{corrige}[Exercice 11 — Expression écrite et orale — type Bac]
\textbf{Tâche écrite} (barème indicatif sur 10 : respect de la tâche et nombre de tours 2 pts ; exactitude grammaticale 3 pts ; lexique de la demande de passeport 2 pts ; formules fonctionnelles 2 pts ; présentation du dialogue 1 pt).

Modèle de dialogue (environ 90 mots) :

\eng{Applicant: Good morning, Sir. I would like to apply for a passport, please.}

\eng{Officer: Good morning. Could you fill in this application form, please?}

\eng{Applicant: Certainly. What documents do I need to provide?}

\eng{Officer: A certified copy of your birth certificate, two recent passport photos and your identity card.}

\eng{Applicant: How much is the fee?}

\eng{Officer: It is 75,000 FCFA, payable at the cashier's desk.}

\eng{Applicant: And how long does it take?}

\eng{Officer: About three weeks. Come back with your receipt.}

\eng{Applicant: Thank you very much for your help.}

\eng{Officer: You're welcome. Have a nice day.}

\textbf{Tâche orale} — réponses personnelles attendues en phrases complètes : \eng{My name is Amina Mbarga.} / \eng{I was born on 12 March 2007 in Bamenda.} / \eng{I am a student at Lycée de Nkol-Eton.} / \eng{I live at Carrefour Emana, Yaoundé.}

Points de vigilance : au moins huit tours de parole et quatre formules fonctionnelles ; alternance correcte demandeur/agent ; \eng{How much} pour le prix, \eng{How long} pour la durée ; à l'oral, soigner l'intonation montante des questions polies (\eng{Could you...?}, \eng{How much...?}) et la date de naissance dite en anglais (\eng{I was born on...}).
\end{corrige}

\newpage

\coursec{Fiche bilan}

\begin{popfigure}
\begin{tikzpicture}[mindmap, grow cyclic, every node/.style={concept, text=white, align=center, font=\small},
root concept/.append style={concept color=popblue, font=\bfseries},
level 1/.append style={level distance=3.4cm, sibling angle=51.4},
level 2/.append style={level distance=2.2cm, sibling angle=45, font=\scriptsize}]
\node[root concept]{Applying for a passport}
child[concept color=poporange]{ node {Key vocabulary}
  child { node {form, photo, fee} }
  child { node {birth certificate} }
}
child[concept color=popgreen]{ node {Asking for information}
  child { node {Could you tell me...?} }
  child { node {How much...? How long...?} }
}
child[concept color=poppurple]{ node {Giving information}
  child { node {You need to...} }
  child { node {It costs..., it takes...} }
}
child[concept color=poppink]{ node {Checking, concluding}
  child { node {Sorry, could you repeat?} }
  child { node {Thank you for your help.} }
}
child[concept color=popgold]{ node {Grammar tools}
  child { node {Wh- questions} }
  child { node {Modals: can, must} }
}
child[concept color=popdark]{ node {Pronunciation}
  child { node {rising intonation} }
  child { node {sentence stress} }
}
child[concept color=popblue!70]{ node {Role play}
  child { node {applicant / officer} }
};
\end{tikzpicture}
\legende{Carte mentale de la leçon : échanger des informations pour demander un passeport.}
\end{popfigure}

\begin{aretenir}[L'essentiel de la leçon]
\textbf{1. Lexique clé (à connaître par cœur)}
\begin{itemize}
\item \cle{passport} : passeport ; \cle{to apply for} : faire une demande de ; \cle{application form} : formulaire de demande
\item \cle{birth certificate} : acte de naissance ; \cle{ID card} : carte d'identité ; \cle{passport-sized photo} : photo d'identité (format passeport)
\item \cle{fee} : frais (officiels) ; \cle{receipt} : reçu ; \cle{processing time} : délai de traitement ; \cle{to collect} : retirer (un document)
\item \cle{immigration office} : bureau de l'immigration ; \cle{requirements} : pièces exigées ; \cle{valid} : valide ; \cle{to expire} : expirer
\end{itemize}

\textbf{2. Formules fonctionnelles et grammaire}

\begin{center}
\begin{tabularx}{\linewidth}{|X|X|}
\hline
\rowcolor{popblueL} Fonction & Formule anglaise \\
\hline
Demander une information & \eng{Could you tell me how to apply for a passport?} / \eng{What documents do I need?} \\
\hline
Demander un prix, un délai & \eng{How much does it cost?} / \eng{How long does it take?} \\
\hline
Donner une information & \eng{You need to fill in this form.} / \eng{It takes about two weeks.} \\
\hline
Relancer, vérifier & \eng{Sorry, could you repeat that, please?} / \eng{Do you mean I must bring two photos?} \\
\hline
Conclure & \eng{Thank you very much for your help.} / \eng{I'll come back next week.} \\
\hline
\end{tabularx}
\end{center}

\begin{itemize}
\item Questions en \cle{Wh-} : mot interrogatif + auxiliaire + sujet + verbe : \eng{Where can I get the form?} (jamais \emph{Where I can get...}).
\item Question indirecte polie : \eng{Could you tell me where the office is?} → ordre sujet + verbe, sans inversion.
\item Modaux : \cle{can/could} (politesse), \cle{must} (obligation), \cle{need to} + base verbale.
\item Prépositions : \eng{apply \textbf{for}}, \eng{fill \textbf{in} a form}, \eng{pay \textbf{for}}, \eng{wait \textbf{for}}.
\end{itemize}

\textbf{3. Méthode express pour le jeu de rôle}
\begin{enumerate}
\item Saluer : \eng{Good morning, sir/madam.}
\item Poser une question claire avec \eng{Could you...?} ou un mot en \emph{Wh-}.
\item Écouter, noter les chiffres (prix, délais) et faire répéter si besoin.
\item Réagir et reformuler : \eng{So I need four photos and my birth certificate.}
\item Remercier et prendre congé.
\end{enumerate}

\textbf{4. Prononciation}
\begin{itemize}
\item Intonation \cle{montante} pour les questions fermées (\eng{Do I need a photo?} ↗) ; \cle{descendante} pour les questions en \emph{Wh-} (\eng{Where is the office?} ↘).
\item Accentuer les mots porteurs de sens : \eng{HOW much DOES it COST?}
\item \eng{passport} se prononce « passe-porte » (a bref) ; \eng{certificate} « se-ti-fi-kète ».
\end{itemize}
\end{aretenir}

\begin{propriete}[Rappel de grammaire : questions directes et indirectes]
\textbf{Question directe} : inversion sujet--auxiliaire obligatoire.

\eng{How long does it take?} « Combien de temps cela prend-il ? »

\eng{What documents do I need?} « Quels documents me faut-il ? »

\textbf{Question indirecte} (après \eng{Could you tell me...}, \eng{Do you know...}, \eng{I'd like to know...}) : on revient à l'ordre \emph{sujet + verbe}, sans auxiliaire \emph{do/does/did} et sans point d'interrogation si la phrase introductive est affirmative.

\eng{Could you tell me where the immigration office is?} « Pourriez-vous me dire où se trouve le bureau de l'immigration ? »

\eng{Do you know how much it costs?} « Savez-vous combien cela coûte ? »

\textbf{Exception} : quand le mot interrogatif est le sujet du verbe, il n'y a jamais d'inversion ni d'auxiliaire : \eng{Who issues passports in Cameroon?} « Qui délivre les passeports au Cameroun ? »
\end{propriete}

\begin{attention}[Erreurs fréquentes des francophones]
\begin{itemize}
\item \emph{« How many it costs? »} → \eng{How much does it cost?} (prix = \emph{how much}, auxiliaire \emph{does}).
\item \emph{« Could you tell me where is the office? »} → \eng{Could you tell me where the office is?} (pas d'inversion dans la question indirecte).
\item \emph{« I need an information. »} → \eng{I need some information.} (\emph{information} est indénombrable en anglais).
\item \emph{« I want to demand a passport. »} → \eng{I would like to apply for a passport.} (\emph{to demand} = exiger ; \emph{to apply for} = faire une demande).
\item \emph{« Fill the form. »} → \eng{Fill in the form.} (phrasal verb : \emph{fill in} un formulaire).
\item \emph{« The passport is valid since ten years. »} → \eng{The passport is valid for ten years.} (durée = \emph{for}).
\end{itemize}
\end{attention}

\begin{exemplebox}[Mini-dialogue modèle à mémoriser]
\eng{— Good morning, madam. I would like to apply for a passport. Could you tell me what documents I need?}

« — Bonjour madame. Je voudrais faire une demande de passeport. Pourriez-vous me dire quels documents il me faut ? »

\eng{— Good morning. You need your birth certificate, your ID card and four passport-sized photos.}

« — Bonjour. Il vous faut votre acte de naissance, votre carte d'identité et quatre photos d'identité. »

\eng{— How much is the fee, and how long does it take?}

« — Quel est le montant des frais, et combien de temps cela prend-il ? »

\eng{— The fee is 75,000 francs and the processing time is about two weeks. You can collect your passport here.}

« — Les frais sont de 75 000 francs et le délai de traitement est d'environ deux semaines. Vous pourrez retirer votre passeport ici. »

\eng{— Thank you very much for your help. Goodbye.}

« — Merci beaucoup pour votre aide. Au revoir. »
\end{exemplebox}

\begin{aretenir}[Checklist avant le Bac]
$\square$ Je connais le lexique du passeport : \eng{application form, fee, receipt, birth certificate, processing time}.

$\square$ Je sais saluer et me présenter poliment : \eng{Good morning, I would like to apply for a passport.}

$\square$ Je sais poser une question directe correcte : \eng{How long does it take?} (inversion avec l'auxiliaire).

$\square$ Je sais poser une question indirecte polie : \eng{Could you tell me where the office is?} (sans inversion).

$\square$ Je sais demander un prix (\eng{How much...?}) et un délai (\eng{How long...?}).

$\square$ Je sais donner une consigne avec \eng{You need to...} et \eng{You must...}.

$\square$ Je sais faire répéter : \eng{Sorry, could you repeat that, please?}

$\square$ Je sais reformuler pour vérifier : \eng{Do you mean that...?} / \eng{So, I need...}

$\square$ Je sais conclure le dialogue : \eng{Thank you for your help. Goodbye.}

$\square$ Je maîtrise l'intonation : montante (questions fermées oui/non), descendante (\emph{Wh-} questions).

$\square$ Je connais les prépositions : \eng{apply for, fill in, pay for, wait for}.

$\square$ J'évite les erreurs de francophone : pas de \emph{« How many it costs? »} mais \eng{How much does it cost?}
\end{aretenir}

\findecours

\end{document}
